\documentclass[11pt,a4paper]{article}
\usepackage[margin=1.8cm]{geometry}
\usepackage{amsmath,amssymb,tabularx,array,needspace,fontspec,cancel,tikz,graphicx}
\setmainfont{FreeSerif}
\newfontfamily\urdufont[Script=Arabic]{FreeSerif}
\TeXXeTstate=1
\newcommand{\urdu}[1]{{\urdufont\beginR #1\endR}}
\newcommand{\lr}[1]{\beginL #1\endL}
\setlength{\parindent}{0pt}
\setlength{\parskip}{0pt}
\renewcommand{\arraystretch}{1.15}
\newcommand{\ut}[1]{\text{{\urdufont\beginR #1\endR}}}
\newcounter{sq}
% \usq{text}: auto-numbered Roman numeral label on the right, Urdu (RTL) text
\newcommand{\usq}[1]{\stepcounter{sq}\par\medskip{\urdufont\raggedleft\beginR\textbf{\lr{\Roman{sq}.}}\ #1\endR\par}}
\newcommand{\fd}[3]{\par\smallskip\noindent\hspace*{7em}\begin{tikzpicture}[>=stealth,scale=0.75]
 \draw (0,0) ellipse (0.9 and 1.5); \draw (3.5,0) ellipse (0.9 and 1.5);
 \foreach \l [count=\i] in {#1} {\node[font=\small] (L\i) at (0,1.25-0.7*\i+0.7) {\l};}
 \foreach \r [count=\i] in {#2} {\node[font=\small] (R\i) at (3.5,1.25-0.7*\i+0.7) {\r};}
 #3
 \end{tikzpicture}\par}
% \tfig{tikz code}: geometry figure
\newcommand{\tfig}[1]{\par\smallskip\noindent\hspace*{7em}{\TeXXeTstate=0 \begin{tikzpicture}[scale=0.8,font=\small]#1\end{tikzpicture}}\par}
% \pfig{file}: figure copied (cropped) from the original paper
\newcommand{\pfig}[2]{\par\smallskip\noindent\hspace*{7em}\includegraphics[width=#2]{#1}\par}
% score-distribution table used in several questions
\newcommand{\dist}{\par\smallskip\noindent\hspace*{7em}\begin{tabular}{|l|c|c|c|c|c|c|c|}\hline
\textbf{Scores} & 24-28 & 29-33 & 34-38 & 39-43 & 44-48 & 49-53 & \textbf{Total}\\ \hline
\textbf{No. of students} & 3 & 6 & 12 & 23 & 15 & 6 & 65\\ \hline\end{tabular}\par\smallskip\hspace*{7em}}
% biology marks tables (6 columns / 5 columns as in the original)
\newcommand{\bioA}{\par\smallskip\noindent\hspace*{7em}\begin{tabular}{|l|c|c|c|c|c|c|}\hline
\textbf{Marks} & 20-24 & 25-29 & 30-34 & 35-39 & 40-44 & 45-49\\ \hline
\textbf{No. of students} & 05 & 08 & 12 & 15 & 03 & 02\\ \hline\end{tabular}\par\smallskip\hspace*{7em}}
\newcommand{\bioB}{\par\smallskip\noindent\hspace*{7em}\begin{tabular}{|l|c|c|c|c|c|}\hline
\textbf{Marks} & 20-24 & 25-29 & 30-34 & 35-39 & 40-44\\ \hline
\textbf{No. of students} & 05 & 08 & 12 & 15 & 03\\ \hline\end{tabular}\par\smallskip\hspace*{7em}}
% continuation paragraph (no label) and Urdu tables
\newcommand{\usc}[1]{\par\smallskip{\urdufont\raggedleft\beginR #1\endR\par}}
\newcommand{\udist}{\par\smallskip\noindent\hspace*{7em}\begin{tabular}{|c|c|c|c|c|c|c|c|}\hline
\urdu{\textbf{نمبرز}} & 24-28 & 29-33 & 34-38 & 39-43 & 44-48 & 49-53 & \urdu{\textbf{مجموعہ}}\\ \hline
\urdu{\textbf{طلبہ کی تعداد}} & 3 & 6 & 12 & 23 & 15 & 6 & \textbf{65}\\ \hline\end{tabular}\par}
\newcommand{\usqd}[2]{\usq{#1}\udist\usc{#2}}
\newcommand{\utabA}{\par\smallskip\noindent\hspace*{7em}\begin{tabular}{|c|c|c|}\hline
\urdu{\textbf{اشیا کی لاگت (ہزاروں میں)}} & \urdu{\textbf{مقدار}} & \urdu{\textbf{اشیا}}\\ \hline
35 & 5 & \urdu{واشنگ مشین}\\ 5 & 3 & \urdu{ہیٹر}\\ 13 & 2 & \urdu{چولہا}\\ 1 & 6 & \urdu{ڈسپنسر}\\ \hline\end{tabular}\par}
\newcommand{\utabB}{\par\smallskip\noindent\hspace*{7em}\begin{tabular}{|c|c|c|c|c|c|}\hline
\urdu{\textbf{کمپیوٹر}} & \urdu{\textbf{اسلامیات}} & \urdu{\textbf{ریاضی}} & \urdu{\textbf{سائنس}} & \urdu{\textbf{انگلش}} & \urdu{\textbf{اردو}}\\ \hline
72 & 85 & 90 & 80 & 65 & 78\\ \hline\end{tabular}\par}
\newcommand{\utabC}{\par\smallskip\noindent\hspace*{7em}\begin{tabular}{|c|c|c|c|}\hline
74 & 78 & 74 & \urdu{نمبرز \lr{(x)}}\\ \hline 5 & 3 & 4 & \urdu{اوزان \lr{(w)}}\\ \hline\end{tabular}\par}
% biology tables (Urdu paper): label column on the right; classes in descending order left-to-right
\newcommand{\ubioB}{\par\smallskip\noindent\hspace*{7em}\begin{tabular}{|c|c|c|c|c|c|}\hline
40-44 & 35-39 & 30-34 & 25-29 & 20-24 & \urdu{\textbf{نمبرز}}\\ \hline
03 & 15 & 12 & 08 & 05 & \urdu{\textbf{طلبہ کی تعداد}}\\ \hline\end{tabular}\par}
\newcommand{\ubioA}{\par\smallskip\noindent\hspace*{7em}\begin{tabular}{|c|c|c|c|c|c|c|}\hline
45-49 & 40-44 & 35-39 & 30-34 & 25-29 & 20-24 & \urdu{\textbf{نمبرز}}\\ \hline
02 & 03 & 15 & 12 & 08 & 05 & \urdu{\textbf{طلبہ کی تعداد}}\\ \hline\end{tabular}\par}
\newcommand{\usqb}[3]{\usq{45 طلبہ کے بائیولوجی کے ماہانہ ٹیسٹ میں درج ذیل نمبر حاصل کیے:}#1\usc{اوپر دی ہوئی تفصیل سے درج ذیل کے بارے میں بتائیں: #2}}

\begin{document}
\begin{center}
  {\Huge\bfseries The Hope Science Academy}\\[3pt]
  {\small Excellence is our standard not goal}\\[2pt]
  {\small 247-E-1, (Opposite LDA School) Johar Town, Lahore. PH: 0300-8194789, 0322-4157001}
\end{center}
\medskip\noindent
\begin{tabularx}{\linewidth}{|X|X|X|X|}
\hline
\textbf{Student Name:} & \textbf{Roll Num:} & \textbf{Class:} 9th & \textbf{Paper Code:} \\ \hline
\textbf{Subject Name:} Mathematics. & \textbf{Time Allowed:} & \textbf{Total Marks:} 2205 & \textbf{Exam Date:} 09-Oct-2026 \\ \hline
\multicolumn{4}{|l|}{\textbf{Exam Syllabus:} FULL BOOK} \\ \hline
\end{tabularx}
\bigskip
\medskip\noindent\textbf{1$\times$831=831}\hfill{\urdufont\beginR\textbf{سوال نمبر \lr{3}۔ مندرجہ ذیل سوالات کے مختصر جوابات تحریر کریں۔}\endR}\par

\usq{درج ذیل کو ناطق یا غیر ناطق اعداد کے طور پر شناخت کریں: \lr{2.353535}}
\usq{درج ذیل کو ناطق یا غیر ناطق اعداد کے طور پر شناخت کریں: \lr{$0.\overline{6}$}}
\usq{درج ذیل کو ناطق یا غیر ناطق اعداد کے طور پر شناخت کریں: \lr{2.236067...}}
\usq{درج ذیل کو ناطق یا غیر ناطق اعداد کے طور پر شناخت کریں: \lr{$\sqrt{7}$}}
\usq{درج ذیل کو ناطق یا غیر ناطق اعداد کے طور پر شناخت کریں: \lr{e}}
\usq{درج ذیل کو ناطق یا غیر ناطق اعداد کے طور پر شناخت کریں: \lr{$\pi$}}
\usq{درج ذیل کو ناطق یا غیر ناطق اعداد کے طور پر شناخت کریں: \lr{$5+\sqrt{11}$}}
\usq{درج ذیل کو ناطق یا غیر ناطق اعداد کے طور پر شناخت کریں: \lr{$\sqrt{3}+\sqrt{13}$}}
\usq{درج ذیل کو ناطق یا غیر ناطق اعداد کے طور پر شناخت کریں: \lr{$\dfrac{15}{4}$}}
\usq{درج ذیل کو ناطق یا غیر ناطق اعداد کے طور پر شناخت کریں: \lr{$(2-\sqrt{2})(2+\sqrt{2})$}}
\usq{درج ذیل کو عددی خط پر ظاہر کریں: \lr{$\sqrt{2}$}}
\usq{درج ذیل کو عددی خط پر ظاہر کریں: \lr{$\sqrt{3}$}}
\usq{درج ذیل کو عددی خط پر ظاہر کریں: \lr{$4\dfrac{1}{3}$}}
\usq{درج ذیل کو عددی خط پر ظاہر کریں: \lr{$-2\dfrac{1}{7}$}}
\usq{درج ذیل کو عددی خط پر ظاہر کریں: \lr{$\dfrac{5}{8}$}}
\usq{درج ذیل کو عددی خط پر ظاہر کریں: \lr{$2\dfrac{3}{4}$}}
\usq{درج ذیل متوالی کسر اعشاریہ کو \lr{$\frac{p}{q}$} کی شکل میں ظاہر کریں جبکہ \lr{p} اور \lr{q} صحیح اعداد ہیں اور \lr{q} ایک غیر صفر عدد ہے: \lr{$0.\overline{4}$}}
\usq{درج ذیل متوالی کسر اعشاریہ کو \lr{$\frac{p}{q}$} کی شکل میں ظاہر کریں جبکہ \lr{p} اور \lr{q} صحیح اعداد ہیں اور \lr{q} ایک غیر صفر عدد ہے: \lr{$0.\overline{37}$}}
\usq{درج ذیل متوالی کسر اعشاریہ کو \lr{$\frac{p}{q}$} کی شکل میں ظاہر کریں جبکہ \lr{p} اور \lr{q} صحیح اعداد ہیں اور \lr{q} ایک غیر صفر عدد ہے: \lr{$0.\overline{21}$}}
\usq{درج ذیل میں استعمال ہونے والی خاصیت کا نام بتائیں: \lr{$(a+4)+b=a+(4+b)$}}
\usq{درج ذیل میں استعمال ہونے والی خاصیت کا نام بتائیں: \lr{$\sqrt{2}+\sqrt{3}=\sqrt{3}+\sqrt{2}$}}
\usq{درج ذیل میں استعمال ہونے والی خاصیت کا نام بتائیں: \lr{$x-x=0$}}
\usq{درج ذیل میں استعمال ہونے والی خاصیت کا نام بتائیں: \lr{$a(b+c)=ab+ac$}}
\usq{درج ذیل میں استعمال ہونے والی خاصیت کا نام بتائیں: \lr{$16+0=16$}}
\usq{درج ذیل میں استعمال ہونے والی خاصیت کا نام بتائیں: \lr{$100\times 1=100$}}
\usq{درج ذیل میں استعمال ہونے والی خاصیت کا نام بتائیں: \lr{$4\times(5\times 8)=(4\times 5)\times 8$}}
\usq{درج ذیل میں استعمال ہونے والی خاصیت کا نام بتائیں: \lr{$ab=ba$}}
\usq{درج ذیل میں استعمال ہونے والی خاصیت کا نام بتائیں: \lr{$-3<-1\Rightarrow 0<2$}}
\usq{درج ذیل میں استعمال ہونے والی خاصیت کا نام بتائیں: \lr{If $a<b$ then $\dfrac{1}{a}>\dfrac{1}{b}$}}
\usq{درج ذیل میں استعمال ہونے والی خاصیت کا نام بتائیں: اگر \lr{$a<b$} ہو تو \lr{$a+c<b+c$}}
\usq{درج ذیل میں استعمال ہونے والی خاصیت کا نام بتائیں: اگر \lr{$ac<bc$} اور \lr{$c>0$} ہو تو \lr{$a<b$}}
\usq{درج ذیل میں استعمال ہونے والی خاصیت کا نام بتائیں: اگر \lr{$ac<bc$} اور \lr{$c<0$} ہو تو \lr{$a>b$}}
\usq{درج ذیل میں استعمال ہونے والی خاصیت کا نام بتائیں: جبکہ \lr{$a>b$} یا \lr{$a=b$} یا \lr{$a<b$}}
\usq{درج ذیل کے درمیان دو ناطق اعداد معلوم کریں: \lr{$\dfrac{1}{3}$} اور \lr{$\dfrac{1}{4}$}}
\usq{درج ذیل کے درمیان دو ناطق اعداد معلوم کریں: \lr{3} اور \lr{4}}
\usq{درج ذیل کے درمیان دو ناطق اعداد معلوم کریں: \lr{$\dfrac{3}{5}$} اور \lr{$\dfrac{4}{5}$}}
\usq{مختتم کسور اعشاریہ کی تعریف کریں۔}
\usq{غیر مختتم کسور اعشاریہ کی تعریف کریں۔}
\usq{ناطق اعداد کی تعریف کریں۔}
\usq{غیر ناطق اعداد کی تعریف کریں۔}
\usq{درج ذیل اعشاری عدد کو ناطق یا غیر ناطق عدد کے طور پر شناخت کریں۔ \lr{0.35}}
\usq{درج ذیل اعشاری عدد کو ناطق یا غیر ناطق عدد کے طور پر شناخت کریں۔ \lr{0.444...}}
\usq{درج ذیل اعشاری عدد کو ناطق یا غیر ناطق عدد کے طور پر شناخت کریں۔ \lr{$3.\overline{5}$}}
\usq{درج ذیل اعشاری عدد کو ناطق یا غیر ناطق عدد کے طور پر شناخت کریں۔ \lr{3.36788542...}}
\usq{درج ذیل اعشاری عدد کو ناطق یا غیر ناطق عدد کے طور پر شناخت کریں۔ \lr{1.709975947...}}
\usq{\lr{$\sqrt{5}$} کو عددی خط پر ظاہر کریں۔}
\usq{درج ذیل متوالی کسر اعشاریہ کو \lr{$\frac{p}{q}$} کی شکل میں ظاہر کریں جبکہ \lr{p} اور \lr{q} صحیح اعداد ہیں: \lr{$0.\overline{5}$}}
\usq{درج ذیل متوالی کسر اعشاریہ کو \lr{$\frac{p}{q}$} کی شکل میں ظاہر کریں جبکہ \lr{p} اور \lr{q} صحیح اعداد ہیں: \lr{$0.\overline{93}$}}
\usq{\lr{2} اور \lr{3} کے درمیان دو ناطق اعداد معلوم کریں۔}
\usq{بیان کو درست ثابت کرنے والی خاصیت کی نشاندہی کریں: اگر \lr{$a>13$} ہو تو \lr{$a+2>15$}}
\usq{بیان کو درست ثابت کرنے والی خاصیت کی نشاندہی کریں: اگر \lr{$3<9$} اور \lr{$6<12$} ہو تو \lr{$9<21$}}
\usq{بیان کو درست ثابت کرنے والی خاصیت کی نشاندہی کریں: اگر \lr{$7>4$} اور \lr{$5>3$} ہو تو \lr{$35>12$}}
\usq{بیان کو درست ثابت کرنے والی خاصیت کی نشاندہی کریں: اگر \lr{$-5<-4\Rightarrow 20>16$}}
\usq{درج ذیل کے مخرج کو ناطق بنائیں: \lr{$\dfrac{13}{4+\sqrt{3}}$}}
\usq{درج ذیل کے مخرج کو ناطق بنائیں: \lr{$\dfrac{\sqrt{2}+\sqrt{5}}{\sqrt{3}}$}}
\usq{درج ذیل کے مخرج کو ناطق بنائیں: \lr{$\dfrac{\sqrt{2}-1}{\sqrt{5}}$}}
\usq{درج ذیل کے مخرج کو ناطق بنائیں: \lr{$\dfrac{6-4\sqrt{2}}{6+4\sqrt{2}}$}}
\usq{درج ذیل کے مخرج کو ناطق بنائیں: \lr{$\dfrac{\sqrt{3}-\sqrt{2}}{\sqrt{3}+\sqrt{2}}$}}
\usq{درج ذیل کے مخرج کو ناطق بنائیں: \lr{$\dfrac{4\sqrt{3}}{\sqrt{7}+\sqrt{5}}$}}
\usq{درج ذیل کو مختصر کریں: \lr{$\left(\dfrac{81}{16}\right)^{-\frac{3}{4}}$}}
\usq{درج ذیل کو مختصر کیجیے: \lr{$\left(\dfrac{3}{4}\right)^{-2}\div\left(\dfrac{4}{9}\right)^{3}\times\dfrac{16}{27}$}}
\usq{درج ذیل کو مختصر کیجیے: \lr{$(0.027)^{-\frac{1}{3}}$}}
\usq{درج ذیل کو مختصر کیجیے: \lr{$\sqrt[7]{\dfrac{x^{14}\times y^{21}\times z^{35}}{y^{14}z^{7}}}$}}
\usq{درج ذیل کو مختصر کیجیے: \lr{$(64)^{-\frac{2}{3}}\div(9)^{-\frac{3}{2}}$}}
\usq{اگر \lr{$x=3+\sqrt{8}$} ہو تو درج ذیل کی قیمت معلوم کریں: \lr{$x+\dfrac{1}{x}$}}
\usq{اگر \lr{$x=3+\sqrt{8}$} ہو تو درج ذیل کی قیمت معلوم کریں: \lr{$x-\dfrac{1}{x}$}}
\usq{اگر \lr{$x=3+\sqrt{8}$} ہو درج ذیل کی قیمت معلوم کریں: \lr{$x^{2}+\dfrac{1}{x^{2}}$}}
\usq{اگر \lr{$x=3+\sqrt{8}$} ہو درج ذیل کی قیمت معلوم کریں: \lr{$\left(x-\dfrac{1}{x}\right)^{2}$}}
\usq{مجذور کی تعریف کریں۔}
\usq{مقادیر اصم کی تعریف کریں۔}
\usq{یک رقمی مقادیر اصم کی تعریف کریں۔}
\usq{دو رقمی مقادیر اصم کی تعریف کریں۔}
\usq{کانجوگیٹ مقادیر اصم کی تعریف کریں۔}
\usq{کیا \lr{0} ایک ناطق نمبر ہے؟ وضاحت کریں۔}
\usq{حقیقی اعداد کی ٹرائکوٹومی پراپرٹی بیان کریں۔}
\usq{\lr{4} اور \lr{5} کے درمیان دو ناطق اعداد معلوم کریں۔}
\usq{اگر \lr{$a=\dfrac{3}{2},b=\dfrac{5}{3}$} اور \lr{$c=\dfrac{7}{5}$} ہو تو ثابت کیجیے: \lr{$a(b+c)=ab+ac$}}
\usq{اگر \lr{$a=\dfrac{3}{2},b=\dfrac{5}{3}$} اور \lr{$c=\dfrac{7}{5}$} ہو تو ثابت کیجیے: \lr{$(a+b)c=ac+bc$}}
\usq{اگر \lr{$a=\dfrac{4}{3},b=\dfrac{5}{2},c=\dfrac{7}{4}$} ہو تو حقیقی اعداد کی خاصیت تلازم بلحاظ جمع اور ضرب کو ثابت کیجیے۔}
\usq{حقیقی اعداد کی ثلاثی خاصیت بیان کریں۔}
\usq{درج ذیل کو مختصر کیجیے: \lr{$\sqrt[5]{\dfrac{x^{15}y^{35}}{z^{20}}}$}}
\usq{درج ذیل کو مختصر کیجیے: \lr{$\sqrt[3]{(27)^{2x}}$}}
\usq{درج ذیل کو مختصر کیجیے: \lr{$\dfrac{6(3)^{n+2}}{3^{n+1}-3^{n}}$}}
\usq{درج ذیل کو مختصر کریں: \lr{$\sqrt[4]{16x^{4}y^{8}}$}}
\usq{درج ذیل کو مختصر کریں: \lr{$\sqrt[3]{27x^{6}y^{9}z^{3}}$}}
\usq{درج ذیل کو مختصر کریں: \lr{$(64)^{-\frac{4}{3}}$}}
\usq{درج ذیل کے مخرج کو ناطق بنائیں: \lr{$\dfrac{3}{\sqrt{5}+\sqrt{2}}$}}
\usq{درج ذیل کے مخرج کو ناطق بنائیں: \lr{$\dfrac{3}{\sqrt{5}-\sqrt{3}}$}}
\usq{تین کیے بعد دیگرے اعداد کا مجموعہ بیالیس ہے، تین اعداد معلوم کریں}
\usq{دو حقیقی نمبروں کا مجموعہ \lr{8} ہے اور ان کا فرق \lr{2} ہے۔ اعداد معلوم کریں۔}
\usq{انسانی جسم کا عمومی درجہ حرارت \lr{98.6 $^{\circ}$F} ہے۔ اسے سیلسیس اور کیلون اسکیل میں تبدیل کریں۔}
\usq{درج ذیل کو سائنسی ترقیم میں لکھیے: \lr{2000000}}
\usq{درج ذیل کو سائنسی ترقیم میں لکھیے: \lr{48900}}
\usq{درج ذیل کو سائنسی ترقیم میں لکھیے: \lr{0.0042}}
\usq{درج ذیل کو سائنسی ترقیم میں لکھیے: \lr{0.0000009}}
\usq{درج ذیل کو سائنسی ترقیم میں لکھیے: \lr{$73\times10^{3}$}}
\usq{درج ذیل کو سائنسی ترقیم میں لکھیے: \lr{$0.65\times10^{2}$}}
\usq{درج ذیل کو عام ترقیم میں لکھیے: \lr{$8.04\times10^{2}$}}
\usq{درج ذیل کو عام ترقیم میں لکھیے: \lr{$3\times10^{5}$}}
\usq{درج ذیل کو عام ترقیم میں لکھیے: \lr{$1.5\times10^{-2}$}}
\usq{درج ذیل کو عام ترقیم میں لکھیے: \lr{$1.77\times10^{7}$}}
\usq{درج ذیل کو عام ترقیم میں لکھیے: \lr{$5.5\times10^{-6}$}}
\usq{درج ذیل کو عام ترقیم میں لکھیے: \lr{$4\times10^{-5}$}}
\usq{روشنی کی رفتار تقریباً \lr{$3\times10^{8}$} میٹر فی سیکنڈ ہے۔ اسے عام ترقیم میں لکھیے۔}
\usq{خط استوا پر زمین کا محیط \lr{40075000} میٹر ہے۔ اس کو سائنسی ترقیم میں لکھیے۔}
\usq{مریخ کا قطر \lr{$6.7779\times10^{3}$} کلو میٹر ہے۔ اس عدد کو عام ترقیم میں لکھیے۔}
\usq{زمین کا قطر \lr{$1.2756\times10^{4}$} کلو میٹر ہے۔ اس عدد کو عام ترقیم میں لکھیے۔}
\usq{\lr{$5.63\times10^{3}$} کو عام ترقیم میں لکھیے۔}
\usq{\lr{$6.6\times10^{-5}$} کو عام ترقیم میں لکھیے۔}
\usq{\lr{29,000,000} کو سائنسی ترقیم میں لکھیے۔}
\usq{\lr{0.000006} کو سائنسی ترقیم میں لکھیے۔}
\usq{سائنسی ترقیم کی تعریف کریں۔}
\usq{سائنسی ترقیم کہاں استعمال ہوتی ہے؟}
\usq{\lr{78,000,000} کو سائنسی ترقیم میں لکھیے۔}
\usq{\lr{0.0000000315} کو سائنسی ترقیم میں لکھیے۔}
\usq{\lr{$3.47\times10^{6}$} کو عام ترقیم میں لکھیے۔}
\usq{\lr{$6.23\times10^{-4}$} کو عام ترقیم میں لکھیے۔}
\usq{درج ذیل کو لوگارتھمی شکل میں لکھیے: \lr{$10^{3}=1000$}}
\usq{درج ذیل کو لوگارتھمی شکل میں لکھیے: \lr{$2^{8}=256$}}
\usq{درج ذیل کو لوگارتھمی شکل میں لکھیے: \lr{$3^{-3}=\dfrac{1}{27}$}}
\usq{درج ذیل کو لوگارتھمی شکل میں لکھیے: \lr{$20^{2}=400$}}
\usq{درج ذیل کو لوگارتھمی شکل میں لکھیے: \lr{$16^{-\frac{1}{4}}=\dfrac{1}{2}$}}
\usq{درج ذیل کو لوگارتھمی شکل میں لکھیے: \lr{$11^{2}=121$}}
\usq{درج ذیل کو لوگارتھمی شکل میں لکھیے: \lr{$p=q^{r}$}}
\usq{درج ذیل کو لوگارتھمی شکل میں لکھیے: \lr{$(32)^{\frac{-1}{5}}=\dfrac{1}{2}$}}
\usq{درج ذیل کو قوت نمائی شکل میں لکھیے: \lr{$\log_{5}125=3$}}
\usq{درج ذیل کو قوت نمائی شکل میں لکھیے: \lr{$\log_{2}16=4$}}
\usq{درج ذیل کو قوت نمائی شکل میں لکھیے: \lr{$\log_{23}1=0$}}
\usq{درج ذیل کو قوت نمائی شکل میں لکھیے: \lr{$\log_{5}5=1$}}
\usq{درج ذیل کو قوت نمائی شکل میں لکھیے: \lr{$\log_{2}\dfrac{1}{8}=-3$}}
\usq{درج ذیل کو قوت نمائی شکل میں لکھیے: \lr{$\dfrac{1}{2}=\log_{9}3$}}
\usq{درج ذیل کو قوت نمائی شکل میں لکھیے: \lr{$5=\log_{10}100000$}}
\usq{درج ذیل کو قوت نمائی شکل میں لکھیے: \lr{$\log_{4}\dfrac{1}{16}=-2$}}
\usq{درج ذیل میں \lr{x} کی قیمت معلوم کریں: \lr{$\log_{x}64=3$}}
\usq{درج ذیل میں \lr{x} کی قیمت معلوم کریں: \lr{$\log_{5}1=x$}}
\usq{درج ذیل میں \lr{x} کی قیمت معلوم کریں: \lr{$\log_{x}8=1$}}
\usq{درج ذیل میں \lr{x} کی قیمت معلوم کریں: \lr{$\log_{10}x=-3$}}
\usq{لوگارتھم کی تعریف کریں۔}
\usq{آپ حقیقی عدد کے لوگارتھم کے بارے میں کیا جانتے ہیں؟}
\usq{\lr{$\log_{2}8=3$} کو قوت نمائی شکل میں تبدیل کیجیے۔}
\usq{\lr{$\log_{10}100=2$} کو قوت نمائی شکل میں کیجیے۔}
\usq{درج ذیل میں \lr{x} کی قیمت معلوم کریں: \lr{$\log_{5}25=x$}}
\usq{درج ذیل میں \lr{x} کی قیمت معلوم کریں: \lr{$\log_{2}x=6$}}
\usq{درج ذیل میں \lr{x} کی قیمت معلوم کریں: \lr{$\log_{5}25=x$}}
\usq{درج ذیل کو لوگارتھمی شکل میں لکھیے: \lr{$3^{4}=81$}}
\usq{درج ذیل کو لوگارتھمی شکل میں لکھیے: \lr{$7^{0}=1$}}
\usq{درج ذیل عدد کا خاصہ لکھیے: \lr{5287}}
\usq{درج ذیل عدد کا خاصہ لکھیے: \lr{59.28}}
\usq{درج ذیل عدد کا خاصہ لکھیے: \lr{0.0567}}
\usq{درج ذیل عدد کا خاصہ لکھیے: \lr{234.7}}
\usq{درج ذیل عدد کا خاصہ لکھیے: \lr{0.000049}}
\usq{درج ذیل عدد کا خاصہ لکھیے: \lr{145000}}
\usq{درج ذیل عدد کا لوگارتھم معلوم کیجیے: \lr{43}}
\usq{درج ذیل عدد کا لوگارتھم معلوم کیجیے: \lr{579}}
\usq{درج ذیل عدد کا لوگارتھم معلوم کیجیے: \lr{1.982}}
\usq{درج ذیل عدد کا لوگارتھم معلوم کیجیے: \lr{0.0876}}
\usq{درج ذیل عدد کا لوگارتھم معلوم کیجیے: \lr{0.047}}
\usq{درج ذیل عدد کا لوگارتھم معلوم کیجیے: \lr{0.000354}}
\usq{\lr{x} کی قیمت معلوم کیجیے: \lr{log x = 0.0065}}
\usq{\lr{x} کی قیمت معلوم کیجیے: \lr{1.192}}
\usq{عام لوگارتھم کی تعریف کریں۔}
\usq{قدرتی لاگ کی تعریف کریں۔}
\usq{خاصہ کی تعریف کریں۔}
\usq{مینٹسا کی تعریف کریں۔}
\usq{خاصہ معلوم کرنے کا طریقہ لکھیں۔}
\usq{اینٹی لوگارتھم کی تعریف کریں۔}
\usq{عام اور قدرتی لوگارتھم میں فرق لکھیں۔}
\usq{کیلکولیٹر کا استعمال کیے بغیر درج ذیل کی قیمت معلوم کیجیے: \lr{$\log_{2}18-\log_{2}9$}}
\usq{کیلکولیٹر کا استعمال کیے بغیر درج ذیل کی قیمت معلوم کیجیے: \lr{$\log_{2}64+\log_{2}2$}}
\usq{کیلکولیٹر کا استعمال کیے بغیر درج ذیل کی قیمت معلوم کیجیے: \lr{$2\log 2+\log 25$}}
\usq{درج ذیل کو واحد لوگارتھم کی شکل میں ظاہر کیجیے: \lr{$\dfrac{1}{2}\log 25+2\log 3$}}
\usq{درج ذیل کو واحد لوگارتھم کی شکل میں ظاہر کیجیے: \lr{$\log 9-\log\dfrac{1}{3}$}}
\usq{درج ذیل کو واحد لوگارتھم کی شکل میں ظاہر کیجیے: \lr{$\log_{5}b^{2}.\log_{a}5^{3}$}}
\usq{درج ذیل کو واحد لوگارتھم کی شکل میں ظاہر کیجیے: \lr{$2\log_{3}x+\log_{3}y$}}
\usq{درج ذیل کو واحد لوگارتھم کی شکل میں ظاہر کیجیے: \lr{$4\log_{5}x-\log_{5}y+\log_{5}z$}}
\usq{درج ذیل کو واحد لوگارتھم کی شکل میں ظاہر کیجیے: \lr{$3\ln a+3\ln b-\ln c$}}
\usq{درج ذیل لوگارتھم کو پھیلا کر لکھیے: \lr{$\log\!\left(\dfrac{11}{5}\right)$}}
\usq{درج ذیل لوگارتھم کو پھیلا کر لکھیے \lr{$\log_{5}\sqrt{8a^{6}}$}}
\usq{درج ذیل لوگارتھم کو پھیلا کر لکھیے \lr{$\ln\!\left(\dfrac{a^{2}b}{c}\right)$}}
\usq{درج ذیل لوگارتھم کو پھیلا کر لکھیے \lr{$\log\!\left(\dfrac{xy}{z}\right)^{\frac{1}{9}}$}}
\usq{درج ذیل لوگارتھم کو پھیلا کر لکھیے \lr{$\ln\sqrt[3]{16x^{3}}$}}
\usq{درج ذیل لوگارتھم کو پھیلا کر لکھیے \lr{$\log_{2}\!\left(\dfrac{1-a}{b}\right)^{5}$}}
\usq{درج ذیل مساوات میں \lr{x} کی قیمت معلوم کیجیے: \lr{$\log 2+\log x=1$}}
\usq{درج ذیل مساوات میں \lr{x} کی قیمت معلوم کیجیے: \lr{$\log_{2}x+\log_{2}8=5$}}
\usq{درج ذیل مساوات میں \lr{x} کی قیمت معلوم کیجیے: \lr{$\log(5x-10)=2$}}
\usq{لوگارتھم کے قوانین لکھیں۔}
\usq{اساس کی تبدیلی کا قانون کیا ہے؟}
\usq{اساس کی تبدیلی کے قانون کو ثابت کریں۔}
\usq{ضرب کے قانون کو ثابت کریں۔}
\usq{تقسیم کے قانون کو ثابت کریں۔}
\usq{قوت نما قانون کو ثابت کریں۔}
\usq{درج ذیل کو لوگارتھم کے قوانین کی مدد سے پھیلا کر لکھیں: \lr{$\log_{3}(20)$}}
\usq{درج ذیل کو لوگارتھم کے قوانین کی مدد سے پھیلا کر لکھیں: \lr{$\log_{2}(9)^{5}$}}
\usq{درج ذیل کو لوگارتھم کے قوانین کی مدد سے پھیلا کر لکھیں: \lr{$\log_{32}27$}}
\usq{درج ذیل کو لوگارتھم کے قوانین کی مدد سے پھیلا کر لکھیے: \lr{$\log_{2}\!\left(\dfrac{x-y}{z}\right)^{3}$}}
\usq{درج ذیل کو لوگارتھم کے قوانین کی مدد سے پھیلا کر لکھیں: \lr{$\log_{5}\!\left(\dfrac{xy}{z}\right)^{8}$}}
\usq{درج ذیل کو واحد لوگارتھم کی شکل میں لکھیے: \lr{$2\log_{3}10-\log_{3}4$}}
\usq{درج ذیل کو واحد لوگارتھم کی شکل میں لکھیے: \lr{$6\log_{3}x+2\log_{3}11$}}
\usq{درج ذیل کو سائنسی ترقیم میں لکھیے: \lr{0.000567}}
\usq{درج ذیل کو سائنسی ترقیم میں لکھیے: \lr{734}}
\usq{درج ذیل کو سائنسی ترقیم میں لکھیے: \lr{$0.33\times10^{3}$}}
\usq{\lr{$2.6\times10^{3}$} درج ذیل کو عام ترقیم میں لکھیے۔}
\usq{درج ذیل کو عام ترقیم میں لکھیے: \lr{$8.794\times10^{-4}$}}
\usq{درج ذیل کو عام ترقیم میں لکھیے: \lr{$6\times10^{-6}$}}
\usq{درج ذیل کو لوگارتھمی شکل میں لکھیے: \lr{$3^{7}=2187$}}
\usq{درج ذیل کو لوگارتھمی شکل میں لکھیے: \lr{$a^{b}=c$}}
\usq{درج ذیل کو لوگارتھمی شکل میں لکھیے: \lr{$(12)^{2}=144$}}
\usq{درج ذیل کو قوت نمائی شکل میں لکھیے: \lr{$\log_{4}8=x$}}
\usq{درج ذیل کو قوت نمائی شکل میں لکھیے: \lr{$\log_{9}729=3$}}
\usq{درج ذیل کو قوت نمائی شکل میں لکھیے: \lr{$\log_{4}1024=5$}}
\usq{درج ذیل میں \lr{x} کی قیمت معلوم کریں: \lr{$\log_{9}x=0.5$}}
\usq{درج ذیل میں \lr{x} کی قیمت معلوم کریں: \lr{$\left(\dfrac{1}{9}\right)^{3x}=27$}}
\usq{درج ذیل کو واحد لوگارتھم کی شکل میں لکھیے: \lr{$7\log x-3\log y^{2}$}}
\usq{رج ذیل کو واحد لوگارتھم کی شکل میں لکھیے: \lr{$3\log 4-\log 32$}}
\usq{درج ذیل کو واحد لوگارتھم کی شکل میں لکھیے: \lr{$\dfrac{1}{3}(\log_{5}8+\log_{5}27)-\log_{5}3$}}
\usq{درج ذیل کو لوگارتھم کے قوانین کی مدد سے پھیلا کر لکھیے: \lr{$\log(x\,y\,z^{6})$}}
\usq{درج ذیل کو لوگارتھم کے قوانین کی مدد سے پھیلا کر لکھیے: \lr{$\log_{3}\sqrt[6]{m^{5}n^{3}}$}}
\usq{درج ذیل کو لوگارتھم کے قوانین کی مدد سے پھیلا کر لکھیے: \lr{$\log\sqrt{8x^{3}}$}}
\usq{درج ذیل کی قیمت لوگارتھم جدول کی مدد سے معلوم کیجیے: \lr{$\sqrt[3]{68.24}$}}
\usq{درج ذیل کو ترقیم سیٹ ساز میں لکھیے: \lr{$\{1,4,9,16,25,36,\ldots,484\}$}}
\usq{درج ذیل کو ترقیم سیٹ ساز میں لکھیے: \lr{$\{2,4,8,16,32,64,\ldots,156\}$}}
\usq{درج ذیل کو ترقیم سیٹ ساز میں لکھیے: \lr{$\{0,\pm1,\pm2,\ldots,\pm1000\}$}}
\usq{درج ذیل کو ترقیم سیٹ ساز میں لکھیے: \lr{$\{6,12,18,\ldots,120\}$}}
\usq{درج ذیل کو ترقیم سیٹ ساز میں لکھیے: \lr{$\{100,102,104,\ldots,400\}$}}
\usq{درج ذیل کو ترقیم سیٹ ساز میں لکھیے: \lr{$\{1,3,9,27,81,\ldots\}$}}
\usq{درج ذیل کو ترقیم سیٹ ساز میں لکھیے: \lr{$\{1,2,4,5,10,20,25,50,100\}$}}
\usq{درج ذیل کو ترقیم سیٹ ساز میں لکھیے: \lr{$\{5,10,15,\ldots,100\}$}}
\usq{درج ذیل کو ترقیم سیٹ ساز میں لکھیے: \lr{$-100$} اور \lr{1000} کے درمیان تمام صحیح اعداد کا سیٹ}
\usq{درج ذیل کو اندراجی طریقہ میں لکھیے: \lr{$\{x|x \text{ is a multiple of } 3\wedge x\leq35\}$}}
\usq{درج ذیل کو اندراجی طریقہ میں لکھیے: \lr{$\{x|x\in R\wedge 2x+1=0\}$}}
\usq{درج ذیل کو اندراجی طریقہ میں لکھیے: \lr{$\{x|x\in P\wedge x<12\}$}}
\usq{درج ذیل کو اندراجی طریقہ میں لکھیے: \lr{$\{x|x \text{ is a divisor of } 128\}$}}
\usq{درج ذیل کو اندراجی طریقہ میں لکھیے: \lr{$\{x|x=2^{n},n\in N\wedge n<8\}$}}
\usq{درج ذیل کو اندراجی طریقہ میں لکھیے: \lr{$\{x|x\in N\wedge x+4=0\}$}}
\usq{درج ذیل کو اندراجی طریقہ میں لکھیے: \lr{$\{x|x\in N\wedge x=x\}$}}
\usq{درج ذیل کو اندراجی طریقہ میں لکھیے: \lr{$\{x|x\in Z\wedge 3x+1=0\}$}}
\usq{درج ذیل سیٹ کے دو واجب تختی سیٹ لکھیں: \lr{$\{a,b,c\}$}}
\usq{درج ذیل سیٹ کے دو واجب تختی سیٹ لکھیں: \lr{$\{0,1\}$}}
\usq{درج ذیل سیٹ کے دو واجب تختی سیٹ لکھیں: \lr{$N$}}
\usq{درج ذیل سیٹ کے دو واجب تختی سیٹ لکھیں: \lr{$Z$}}
\usq{درج ذیل سیٹ کے دو واجب تختی سیٹ لکھیں: \lr{$Q$}}
\usq{درج ذیل سیٹ کے دو واجب تختی سیٹ لکھیں: \lr{$R$}}
\usq{درج ذیل سیٹ کے دو واجب تختی سیٹ لکھیں: \lr{$\{x|x\in Q\wedge 0<x\leq2\}$}}
\usq{کیا کوئی ایسا سیٹ ہوتا ہے جس کا کوئی واجب تختی سیٹ نہیں ہوتا؟ اگر ایسا ہے تو اس سیٹ کا نام لکھیے۔}
\usq{\lr{$\{a,b\}$} اور \lr{$\{\{a,b\}\}$} میں کیا فرق ہے؟}
\usq{درج ذیل سیٹ کے پاور سیٹ کے ارکان کی تعداد کیا ہے؟ \lr{$\{\ \}$}}
\usq{درج ذیل سیٹ کے پاور سیٹ کے ارکان کی تعداد کیا ہے؟ \lr{$\{0,1\}$}}
\usq{درج ذیل سیٹ کے پاور سیٹ کے ارکان کی تعداد کیا ہے؟ \lr{$\{1,2,3,4,5,6,7\}$}}
\usq{درج ذیل سیٹ کے پاور سیٹ کے ارکان کی تعداد کیا ہے؟ \lr{$\{0,1,2,3,4,5,6,7\}$}}
\usq{درج ذیل سیٹ کے پاور سیٹ کے ارکان کی تعداد کیا ہے؟ \lr{$\{a,\{b,c\}\}$}}
\usq{درج ذیل سیٹ کے پاور سیٹ کے ارکان کی تعداد کیا ہے؟ \lr{$\{\{a,b\},\{b,c\},\{d,e\}\}$}}
\usq{درج ذیل سیٹ کا پاور سیٹ لکھیں: \lr{$\{9,11\}$}}
\usq{درج ذیل سیٹ کا پاور سیٹ لکھیں: \lr{$\{+,-,\times,\div\}$}}
\usq{درج ذیل سیٹ کا پاور سیٹ لکھیں: \lr{$\{\phi\}$}}
\usq{درج ذیل سیٹ کا پاور سیٹ لکھیں: \lr{$\{a,\{b,c\}\}$}}
\usq{سیٹ کی تعریف کریں۔}
\usq{بیانیہ طریقہ کیا ہے؟}
\usq{اندراجی طریقہ کیا ہے؟}
\usq{ترقیم سیٹ ساز طریقہ کیا ہے؟}
\usq{یک رکنی سیٹ کیا ہے؟}
\usq{نل سیٹ کی تعریف کریں۔}
\usq{خالی سیٹ کو کیسے لکھا جاتا ہے؟}
\usq{مساوی سیٹ کی تعریف کریں۔}
\usq{مترادف سیٹ کی تعریف کریں۔}
\usq{تختی سیٹ کی تعریف کریں۔}
\usq{واجب تختی سیٹ کی تعریف کریں۔}
\usq{غیر واجب تختی سیٹ کی تعریف کریں۔}
\usq{یونیورسل سیٹ کی تعریف کریں۔}
\usq{پاور سیٹ کی تعریف کریں۔}
\usq{یونیورسل سیٹ پر غور کریں: \lr{$U=\{x|\ x\ \ut{\lr{2} کا ضعف ہے}\ \wedge\ 0<x\leq30\}$}، \lr{$A=\{x|\ x\ \ut{\lr{6} کا ضعف ہے}\}$} اور \lr{$B=\{x|\ x\ \ut{\lr{8} کا ضعف ہے}\}$}۔ سیٹ \lr{A} اور \lr{B} کے تمام ارکان کو اندراجی طریقہ میں لکھیں۔}
\usq{یونیورسل سیٹ پر غور کریں: \lr{$U=\{x|\ x\ \ut{\lr{2} کا ضعف ہے}\ \wedge\ 0<x\leq30\}$}، \lr{$A=\{x|\ x\ \ut{\lr{6} کا ضعف ہے}\}$} اور \lr{$B=\{x|\ x\ \ut{\lr{8} کا ضعف ہے}\}$} ہو تو \lr{$A\cap B$} معلوم کریں۔}
\usq{اگر \lr{$U=\{1,2,3,\ldots,20\}$} اور \lr{$A=\{1,3,5,\ldots,19\}$} ہو تو درج ذیل ثابت کریں: \lr{$A\cap U=A$}}
\usq{دو سیٹوں کی یونین کیا ہے؟}
\usq{دو سیٹوں کا تقاطع کیا ہے؟}
\usq{غیر مشترک سیٹ کی تعریف کریں۔}
\usq{متراکب سیٹ کی تعریف کریں۔}
\usq{سیٹ کا کمپلیمنٹ کیا ہے؟}
\usq{وین اشکال کیا ہے؟}
\usq{وین اشکال کو سب سے پہلے کس نے استعمال کیا؟}
\usq{شمولیت واخراج کا اصول کیا ہے؟}
\usq{کارڈنیلٹی کی تعریف کریں۔}
\usq{\lr{$A=\{1,2,3,4\}$} کے لیے درج ذیل ربط \lr{A} میں معلوم کریں۔ ہر ربط کی ڈومین اور رینج بھی معلوم کریں۔ \lr{$\{(x,y)\mid y=x\}$}}
\usq{\lr{$A=\{1,2,3,4\}$} کے لیے درج ذیل ربط \lr{A} میں معلوم کریں۔ ہر ربط کی ڈومین اور رینج بھی معلوم کریں۔ \lr{$\{(x,y)\mid y+x=5\}$}}
\usq{\lr{$A=\{1,2,3,4\}$} کے لیے درج ذیل ربط \lr{A} میں معلوم کریں۔ ہر ربط کی ڈومین اور رینج بھی معلوم کریں۔ \lr{$\{(x,y)\mid x+y<5\}$}}
\usq{\lr{$A=\{1,2,3,4\}$} کے لیے درج ذیل ربط \lr{A} میں معلوم کریں۔ ہر ربط کی ڈومین اور رینج بھی معلوم کریں۔ \lr{$\{(x,y)\mid x+y>5\}$}}
\usq{دی گئی شکل کون سی تفاعل کو ظاہر کرتی ہے اور کس قسم کا تفاعل ہے؟}
\fd{1,2,3}{a,b,c,d}{\draw[->](L1)--(R1); \draw[->](L1)--(R2); \draw[<->](L2)--(R3); \draw[<->](L3)--(R4);}
\usq{دی گئی شکل کون سی تفاعل کو ظاہر کرتی ہے اور کس قسم کا تفاعل ہے؟}
\fd{a,b,c}{1,3,5}{\draw[->](L1)--(R1); \draw[->](L2)--(R2); \draw[->](L3)--(R3);}
\usq{دی گئی شکل کون سی تفاعل کو ظاہر کرتی ہے اور کس قسم کا تفاعل ہے؟}
\fd{1,2,3}{a,b,c}{\draw[->](L1)--(R1); \draw[->](L2)--(R2); \draw[->](L3)--(R3);}
\usq{دی گئی شکل کون سی تفاعل کو ظاہر کرتی ہے اور کس قسم کا تفاعل ہے؟}
\fd{l,m,n}{x,y,z}{\draw[->](L1)--(R1); \draw[->](L2)--(R1); \draw[->](L3)--(R3);}
\usq{اگر \lr{$g(x)=3x+2$} ہو تو درج ذیل معلوم کریں۔ \lr{g(0)}}
\usq{اگر \lr{$g(x)=3x+2$} ہو تو درج ذیل معلوم کریں۔ \lr{$g(-3)$}}
\usq{اگر \lr{$g(x)=3x+2$} ہو تو درج ذیل معلوم کریں۔ \lr{$g\!\left(\dfrac{2}{3}\right)$}}
\usq{اگر \lr{$h(x)=x^{2}+1$} ہو تو درج ذیل معلوم کریں۔ \lr{$h(1)$}}
\usq{اگر \lr{$h(x)=x^{2}+1$} ہو تو درج ذیل معلوم کریں۔ \lr{$h(-4)$}}
\usq{اگر \lr{$h(x)=x^{2}+1$} ہو تو درج ذیل معلوم کریں۔ \lr{$h\!\left(-\dfrac{1}{2}\right)$}}
\usq{فرض کریں۔ \lr{$f(x)=5x+1$} اگر \lr{$f(x)=32$} تو \lr{x} کی قیمت معلوم کریں۔}
\usq{تفاعل کی تعریف کریں۔}
\usq{تفاعل کی کتنی اقسام ہیں؟ ان کے نام لکھیں۔}
\usq{انٹو تفاعل کیا ہے؟}
\usq{آن ٹو تفاعل کیا ہے؟}
\usq{ون ون تفاعل کیا ہے؟}
\usq{بائی جیکٹیو تفاعل کیا ہے؟}
\usq{خطی تفاعل کی تعریف کریں۔}
\usq{دو درجی تفاعل کی تعریف کریں۔}
\usq{ثنائی ربط کیا ہے؟}
\usq{ڈومین کی تعریف کریں۔}
\usq{رینج کی تعریف کریں۔}
\usq{درج ذیل کو اندراجی طریقے میں لکھیں: \lr{$\{x\mid x=2n,\ n\in N\}$}}
\usq{درج ذیل کو اندراجی طریقے میں لکھیں: \lr{$\{x\mid x=2m+1,\ m\in N\}$}}
\usq{درج ذیل کو اندراجی طریقے میں لکھیں: \lr{$\{x\mid x=11n,n\in W\wedge n<11\}$}}
\usq{درج ذیل کو اندراجی طریقے میں لکھیں: \lr{$\{x\mid x\in E\wedge 4<x<6\}$}}
\usq{درج ذیل کو اندراجی طریقے میں لکھیں: \lr{$\{x\mid x\in O\wedge 5\leq x<7\}$}}
\usq{درج ذیل کو اندراجی طریقے میں لکھیں: \lr{$\{x\mid x\in Q\wedge x^{2}=2\}$}}
\usq{درج ذیل کو اندراجی طریقے میں لکھیں: \lr{$\{x\mid x\in Q\wedge x=-x\}$}}
\usq{درج ذیل کو اندراجی طریقے میں لکھیں: \lr{$\{x\mid x\in R\wedge x\notin Q'\}$}}
\usq{فرض کریں۔ \lr{$U=\{1,2,3,4,5,6,7,8,9,10\}$}، \lr{$A=\{2,4,6,8,10\}$}، \lr{$B=\{1,2,3,4,5\}$} اور \lr{$C=\{1,3,5,7,9\}$} ہو تو درج ذیل سیٹ کے ارکان کی فہرست بنائیں۔ \lr{$A'$}}
\usq{فرض کریں۔ \lr{$U=\{1,2,3,4,5,6,7,8,9,10\}$}، \lr{$A=\{2,4,6,8,10\}$}، \lr{$B=\{1,2,3,4,5\}$} اور \lr{$C=\{1,3,5,7,9\}$} ہو تو درج ذیل سیٹ کے ارکان کی فہرست بنائیں۔ \lr{$B'$}}
\usq{فرض کریں۔ \lr{$U=\{1,2,3,4,5,6,7,8,9,10\}$}، \lr{$A=\{2,4,6,8,10\}$}، \lr{$B=\{1,2,3,4,5\}$} اور \lr{$C=\{1,3,5,7,9\}$} ہو تو درج ذیل سیٹ کے ارکان کی فہرست بنائیں۔ \lr{$A\cup B$}}
\usq{فرض کریں۔ \lr{$U=\{1,2,3,4,5,6,7,8,9,10\}$}، \lr{$A=\{2,4,6,8,10\}$}، \lr{$B=\{1,2,3,4,5\}$} اور \lr{$C=\{1,3,5,7,9\}$} ہو تو درج ذیل سیٹ کے ارکان کی فہرست بنائیں۔ \lr{$A-B$}}
\usq{فرض کریں۔ \lr{$U=\{1,2,3,4,5,6,7,8,9,10\}$}، \lr{$A=\{2,4,6,8,10\}$}، \lr{$B=\{1,2,3,4,5\}$} اور \lr{$C=\{1,3,5,7,9\}$} ہو تو درج ذیل سیٹ کے ارکان کی فہرست بنائیں۔ \lr{$A\cap C$}}
\usq{فرض کریں۔ \lr{$U=\{1,2,3,4,5,6,7,8,9,10\}$}، \lr{$A=\{2,4,6,8,10\}$}، \lr{$B=\{1,2,3,4,5\}$} اور \lr{$C=\{1,3,5,7,9\}$} ہو تو درج ذیل سیٹ کے ارکان کی فہرست بنائیں۔ \lr{$U'$}}
\usq{وین اشکال کا استعمال کرتے ہوئے اگر ضروری ہو تو درج ذیل کے برابر ایک رکنی سیٹ معلوم کریں: \lr{$A'$}}
\usq{وین اشکال کا استعمال کرتے ہوئے اگر ضروری ہو تو درج ذیل کے برابر ایک رکنی سیٹ معلوم کریں: \lr{$A\cap U$}}
\usq{وین اشکال کا استعمال کرتے ہوئے اگر ضروری ہو تو درج ذیل کے برابر ایک رکنی سیٹ معلوم کریں: \lr{$A\cup U$}}
\usq{وین اشکال کا استعمال کرتے ہوئے اگر ضروری ہو تو درج ذیل کے برابر ایک رکنی سیٹ معلوم کریں: \lr{$A\cup\phi$}}
\usq{وین اشکال کا استعمال کرتے ہوئے اگر ضروری ہو تو درج ذیل کے برابر ایک رکنی سیٹ معلوم کریں: \lr{$\phi\cup\phi$}}
\usq{فرض کریں \lr{$P=\{x\mid x=5m,m\in N\}$} اور \lr{$Q=\{x\mid x=2m,m\in N\}$} ہو تو \lr{$P\cap Q$} معلوم کریں۔}
\usq{اگر \lr{$g(x)=7x-2$} ہو تو معلوم کریں: \lr{$g(0)$}}
\usq{اگر \lr{$g(x)=7x-2$} ہو تو معلوم کریں: \lr{$g(-1)$}}
\usq{اگر \lr{$g(x)=7x-2$} ہو تو معلوم کریں: \lr{$g\!\left(-\dfrac{5}{3}\right)$}}
\usq{اگر \lr{$s(x)=8x^{2}-3$} ہو تو معلوم کریں: \lr{$s(1)$}}
\usq{اگر \lr{$s(x)=8x^{2}-3$} ہو تو معلوم کریں: \lr{$s(-9)$}}
\usq{اگر \lr{$s(x)=8x^{2}-3$} ہو تو معلوم کریں: \lr{$s\!\left(\dfrac{7}{2}\right)$}}
\usq{فرض کریں۔ \lr{$k(x)=7x-5$} اگر \lr{$k(x)=100$} ہو تو \lr{x} ی قیمت معلوم کریں۔}
\usq{مشترک جزو ضربی کی شناخت کرکے تجزی کریں: \lr{$6x+12$}}
\usq{مشترک جزو ضربی کی شناخت کرکے تجزی کریں: \lr{$15y^{2}+20y$}}
\usq{مشترک جزو ضربی کی شناخت کرکے تجزی کریں: \lr{$-12x^{2}-3x$}}
\usq{مشترک جزو ضربی کی شناخت کرکے تجزی کریں: \lr{$4a^{2}b+8ab^{2}$}}
\usq{مشترک جزو ضربی کی شناخت کرکے تجزی کریں: \lr{$xy-3x^{2}+2x$}}
\usq{مشترک جزو ضربی کی شناخت کرکے تجزی کریں: \lr{$3a^{2}b-9ab^{2}+15ab$}}
\usq{تجزی کریں: \lr{$5x+15$}}
\usq{تجزی کریں: \lr{$x^{2}+4x+3$}}
\usq{تجزی کریں: \lr{$x^{2}+6x+8$}}
\usq{تجزی کریں: \lr{$x^{2}+4x+4$}}
\usq{تجزی کریں: \lr{$x^{2}+x-12$}}
\usq{تجزی کریں: \lr{$x^{2}+7x+10$}}
\usq{تجزی کریں: \lr{$x^{2}-6x+8$}}
\usq{تجزی کریں: \lr{$x^{2}-x-56$}}
\usq{تجزی کریں: \lr{$x^{2}-10x-24$}}
\usq{تجزی کریں: \lr{$y^{2}+4y-12$}}
\usq{تجزی کریں: \lr{$y^{2}+13y+36$}}
\usq{تجزی کریں: \lr{$x^{2}-x-2$}}
\usq{تجزی کریں: \lr{$2x^{2}+7x+3$}}
\usq{تجزی کریں: \lr{$4x^{2}+13x+3$}}
\usq{تجزی کریں: \lr{$3x^{2}+5x+2$}}
\usq{مشترک جزو ضربی کی تعریف کریں۔}
\usq{دو درجی جملہ کی تعریف کریں۔}
\usq{\lr{$x^{2}+2x$} کے مشترک جزو ضربی ٹھوس، تصویری اور علامتی طریقے سے معلوم کریں۔}
\usq{\lr{$x^{2}-5x+4$} کی تجزی ٹھوس، تصویری اور علامتی طریقے سے کریں۔}
\usq{\lr{$x^{2}-3x-10$} کی تجزی ٹھوس، تصویری اور علامتی طریقے سے کریں۔}
\usq{تجزی کریں: \lr{$x^{2}+9x+14$}}
\usq{تجزی کریں: \lr{$x^{2}-11x+24$}}
\usq{تجزی کریں: \lr{$p^{2}+11p+18$}}
\usq{تجزی کریں: \lr{$3x^{2}-4x-4$}}
\usq{تجزی کریں: \lr{$8x^{3}+12x^{2}+6x+1$}}
\usq{تجزی کریں: \lr{$27a^{3}+108a^{2}b+144ab^{2}+64b^{3}$}}
\usq{تجزی کریں: \lr{$x^{3}+48x^{2}y+108xy^{2}+216y^{3}$}}
\usq{تجزی کریں: \lr{$8x^{3}-125y^{3}+150xy^{2}-60x^{2}y$}}
\usq{تجزی کریں: \lr{$125a^{3}-1$}}
\usq{تجزی کریں: \lr{$64x^{3}+125$}}
\usq{تجزی کریں: \lr{$x^{6}-27$}}
\usq{تجزی کریں: \lr{$1000a^{3}+1$}}
\usq{تجزی کریں: \lr{$343x^{3}+216$}}
\usq{تجزی کریں: \lr{$27-512y^{3}$}}
\usq{تجزی کریں: \lr{$x^{3}-18x^{2}+108x-216$}}
\usq{تجزی کریں: \lr{$8x^{3}+27$}}
\usq{تجزی کریں: \lr{$x^{3}-27y^{3}$}}
\usq{عاداعظم بذریعہ تجزی معلوم کریں: \lr{$21x^{2}y,\ 35xy^{2}$}}
\usq{عاداعظم بذریعہ تجزی معلوم کریں: \lr{$4x^{2}-9y^{2},\ 2x^{2}-3xy$}}
\usq{عاداعظم بذریعہ تجزی معلوم کریں: \lr{$x^{3}-1,\ x^{2}+x+1$}}
\usq{درج ذیل جملے کا ذواضعاف اقل بذریعہ تجزی معلوم کریں: \lr{$2a^{2}b,\ 4ab^{2},\ 6ab$}}
\usq{درج ذیل جملے کا ذواضعاف اقل بذریعہ تجزی معلوم کریں: \lr{$x^{2}+x,\ x^{3}+x^{2}$}}
\usq{درج ذیل جملے کا ذواضعاف اقل بذریعہ تجزی معلوم کریں: \lr{$a^{2}-4a+4,\ a^{2}-2a$}}
\usq{عاداعظم کی تعریف کریں۔}
\usq{ذواضعاف اقل کی تعریف کریں۔}
\usq{عاداعظم اور ذواضعاف اقل کے درمیان تعلق لکھیں۔}
\usq{\lr{$6x^{2}y,\ 9xy^{2}$} کا عاداعظم معلوم کریں۔}
\usq{\lr{$4x^{2}y,\ 8x^{3}y^{2}$} کا ذواضعاف اقل معلوم کریں۔}
\usq{درج ذیل جملے کا جذر المربع بذریعہ تجزی معلوم کریں۔ \lr{$x^{2}-8x+16$}}
\usq{درج ذیل جملے کا جذر المربع بذریعہ تجزی معلوم کریں۔ \lr{$9x^{2}+12x+4$}}
\usq{درج ذیل جملے کا جذر المربع بذریعہ تجزی معلوم کریں۔ \lr{$36a^{2}+84a+49$}}
\usq{درج ذیل جملے کا جذر المربع بذریعہ تجزی معلوم کریں۔ \lr{$64y^{2}-32y+4$}}
\usq{درج ذیل جملے کی تجزی کریں: \lr{$4x^{3}+18x^{2}-12x$}}
\usq{درج ذیل جملے کی تجزی کریں: \lr{$x^{3}+64y^{3}$}}
\usq{درج ذیل جملے کی تجزی کریں: \lr{$-x^{2}-23x-60$}}
\usq{درج ذیل جملے کی تجزی کریں: \lr{$x^{3}y^{3}-8$}}
\usq{درج ذیل جملے کی تجزی کریں: \lr{$2x^{2}+7x+3$}}
\usq{درج ذیل جملے کی تجزی کریں: \lr{$x^{4}+64$}}
\usq{ذواضعاف اقل اور عاداعظم معلوم کریں: \lr{$x^{2}+8x+16,\ x^{2}-16$}}
\usq{ذواضعاف اقل اور عاداعظم معلوم کریں: \lr{$x^{3}-9x,\ x^{2}-4x+3$}}
\usq{\lr{$16x^{4}+8x^{2}+1$} کا جذر المربع بذریعہ تجزی اور بذریعہ تقسیم معلوم کریں}
\usq{درج ذیل جملے کی تجزی کریں: \lr{$x^{4}+2x^{2}+9$}}
\usq{درج ذیل مساوات کو حل کریں اور عددی خط پر ظاہر کریں: \lr{$12x+30=-6$}}
\usq{درج ذیل مساوات کو حل کریں اور عددی خط پر ظاہر کریں: \lr{$\dfrac{x}{3}+6=-12$}}
\usq{درج ذیل مساوات کو حل کریں اور عددی خط پر ظاہر کریں: \lr{$\dfrac{x}{2}-\dfrac{3x}{4}=\dfrac{1}{12}$}}
\usq{درج ذیل مساوات کو حل کریں اور عددی خط پر ظاہر کریں: \lr{$2=7(2x+4)+12x$}}
\usq{درج ذیل غیر مساوات کو حل کریں اور عددی خط پر ظاہر کریں: \lr{$x-6\leq-2$}}
\usq{درج ذیل غیر مساوات کو حل کریں اور عددی خط پر ظاہر کریں: \lr{$-9>-16+x$}}
\usq{درج ذیل غیر مساوات کو حل کریں اور عددی خط پر ظاہر کریں: \lr{$3+2x\geq3$}}
\usq{درج ذیل غیر مساوات کو حل کریں اور عددی خط پر ظاہر کریں: \lr{$6(x+10)\leq0$}}
\usq{درج ذیل یک درجی غیر مساوات کے حل کا خطہ \lr{xy}-مستوی میں سایہ دار کر کے ظاہر کریں: \lr{$2x+y\leq6$}}
\usq{درج ذیل یک درجی غیر مساوات کے حل کا خطہ \lr{xy}-مستوی میں سایہ دار کر کے ظاہر کریں: \lr{$3x+7y\geq21$}}
\usq{درج ذیل یک درجی غیر مساوات کے حل کا خطہ \lr{xy}-مستوی میں سایہ دار کر کے ظاہر کریں: \lr{$3x-2y\geq6$}}
\usq{درج ذیل یک درجی غیر مساوات کے حل کا خطہ \lr{xy}-مستوی میں سایہ دار کر کے ظاہر کریں: \lr{$5x-4y\leq20$}}
\usq{درج ذیل یک درجی غیر مساوات کے حل کا خطہ \lr{xy}-مستوی میں سایہ دار کر کے ظاہر کریں: \lr{$2x+1\geq0$}}
\usq{درج ذیل یک درجی غیر مساوات کے حل کا خطہ \lr{xy}-مستوی میں سایہ دار کر کے ظاہر کریں: \lr{$3y-4\leq0$}}
\usq{یک درجی مساوات کی تعریف کریں۔}
\usq{ایک متغیر میں یک درجی مساوات کو کیسے حل کیا جاتا ہے؟}
\usq{درج ذیل مساوات کو حل کریں اور عددی خط پر ظاہر کریں: \lr{$3x-5=7$}}
\usq{درج ذیل مساوات کو حل کریں اور عددی خط پر ظاہر کریں: \lr{$\dfrac{x-2}{5}-\dfrac{x-4}{2}=2$}}
\usq{\lr{$\dfrac{2}{3}x-1<0$} کو حل کریں اور عددی خط پر ظاہر کریں۔}
\usq{\lr{$x+2y<6$} کو حل کریں۔}
\usq{\lr{xy}-مستوی میں درج ذیل یک درجی غیر مساوات کو حل کریں۔ \lr{$2x\geq-3$}}
\usq{\lr{xy}-مستوی میں درج ذیل یک درجی غیر مساوات کو حل کریں۔ \lr{$y\leq2$}}
\usq{مسئلے کی رکاوٹوں کی تعریف کریں۔}
\usq{تصفیہ کرنے والے متغیرات کی تعریف کریں۔}
\usq{قابل عمل خطے کی تعریف کریں۔}
\usq{قابل عمل حل کی تعریف کریں۔}
\usq{غیر منفی رکاوٹوں کی تعریف کریں۔}
\usq{کونہ نقطہ کی تعریف کریں۔}
\usq{راس کی تعریف کریں۔}
\usq{آبجیکٹو تفاعل کی تعریف کریں۔}
\usq{موافق حل کی تعریف کریں۔}
\usq{حل کریں اور عددی خط پر ظاہر کریں: \lr{$\dfrac{x+5}{3}=1-x$}}
\usq{حل کریں اور عددی خط پر ظاہر کریں: \lr{$\dfrac{2x+1}{3}+\dfrac{1}{2}=1-\dfrac{x-1}{3}$}}
\usq{حل کریں اور عددی خط پر ظاہر کریں: \lr{$3x+7<16$}}
\usq{درج ذیل زاویہ کون سے ربع میں واقع ہے اور اس کا کوٹرمنل زاویہ لکھیں: \lr{$65^{\circ}$}}
\usq{درج ذیل زاویہ کون سے ربع میں واقع ہے اور اس کا کوٹرمنل زاویہ لکھیں: \lr{$135^{\circ}$}}
\usq{درج ذیل زاویہ کون سے ربع میں واقع ہے اور اس کا کوٹرمنل زاویہ لکھیں: \lr{$-40^{\circ}$}}
\usq{درج ذیل زاویہ کون سے ربع میں واقع ہے اور اس کا کوٹرمنل زاویہ لکھیں: \lr{$210^{\circ}$}}
\usq{درج ذیل زاویہ کون سے ربع میں واقع ہے اور اس کا کوٹرمنل زاویہ لکھیں: \lr{$-150^{\circ}$}}
\usq{درج ذیل کو ڈگری، منٹ اور سیکنڈ میں تبدیل کریں: \lr{$123.456^{\circ}$}}
\usq{درج ذیل کو ڈگری، منٹ اور سیکنڈ میں تبدیل کریں: \lr{$58.7891^{\circ}$}}
\usq{درج ذیل کو ڈگری، منٹ اور سیکنڈ میں تبدیل کریں: \lr{$90.5678^{\circ}$}}
\usq{درج ذیل کو اعشاری ڈگری میں تبدیل کریں: \lr{$65^{\circ}32'15''$}}
\usq{درج ذیل کو اعشاری ڈگری میں تبدیل کریں: \lr{$42^{\circ}18'45''$}}
\usq{درج ذیل کو اعشاری ڈگری میں تبدیل کریں: \lr{$78^{\circ}45'36''$}}
\usq{درج ذیل کو ریڈین میں تبدیل کریں: \lr{$36^{\circ}$}}
\usq{درج ذیل کو ریڈین میں تبدیل کریں: \lr{$22.5^{\circ}$}}
\usq{درج ذیل کو ریڈین میں تبدیل کریں: \lr{$67.5^{\circ}$}}
\usq{درج ذیل کو ڈگری میں تبدیل کریں: \lr{$\dfrac{\pi}{16}\,rad$}}
\usq{درج ذیل کو ڈگری میں تبدیل کریں: \lr{$\dfrac{11\pi}{5}\,rad$}}
\usq{درج ذیل کو ڈگری میں تبدیل کریں: \lr{$\dfrac{7\pi}{6}\,rad$}}
\usq{قوس کی لمبائی اور علاقہ کا رقبہ معلوم کریں اگر: \lr{$r=6$ cm} اور مرکزی زاویہ \lr{$\theta=\dfrac{\pi}{3}$} ریڈین}
\usq{قوس کی لمبائی اور علاقہ کا رقبہ معلوم کریں اگر: \lr{$r=\dfrac{4.8}{\pi}$} اور مرکزی زاویہ \lr{$\theta=\dfrac{5\pi}{6}$}}
\usq{اس علاقے کے رقبے کی فی صد معلوم کریں جس کا مرکزی زاویہ \lr{$\dfrac{\pi}{8}$} ریڈین ہو۔}
\usq{تکونیات کیا ہے؟}
\usq{کب زاویہ مثبت ہوگا؟}
\usq{کب زاویہ منفی ہوگا؟}
\usq{کوٹرمینل کیا ہوتے ہیں؟}
\usq{ڈگری کی تعریف کریں۔}
\usq{ریڈین کی تعریف کریں۔}
\usq{مکمل دائرہ کی تعریف کریں۔}
\usq{نصف دائرہ کی تعریف کریں۔}
\usq{قائمہ زاویہ کی تعریف کریں۔}
\usq{ثابت کریں: \lr{$l=r\theta$}}
\usq{ثابت کریں: \lr{$A=\dfrac{1}{2}r^{2}\theta$}}
\usq{\lr{$73.12^{\circ}$} کو ڈگری، منٹ اور سیکنڈ میں تبدیل کریں۔}
\usq{\lr{$109.42^{\circ}$} کو ڈگری، منٹ اور سیکنڈ میں تبدیل کریں۔}
\usq{\lr{$45^{\circ}45'45''$} کو اعشاری ڈگری میں تبدیل کریں}
\usq{\lr{$94^{\circ}27'54''$} کو اعشاری ڈگری میں تبدیل کریں۔}
\usq{ریڈین سے ڈگری میں تبدیل کریں: \lr{$\dfrac{5\pi}{3}\,rad$}}
\usq{ریڈین سے ڈگری میں تبدیل کریں: \lr{$\dfrac{7\pi}{6}\,rad$}}
\usq{ریڈین سے ڈگری میں تبدیل کریں: \lr{$\dfrac{11\pi}{6}\,rad$}}
\usq{ریڈین سے ڈگری میں تبدیل کریں: \lr{1.2rad}}
\usq{ڈگری سے ریڈین میں تبدیل کریں: \lr{$15^{\circ}$}}
\usq{ڈگری سے ریڈین میں تبدیل کریں: \lr{$75^{\circ}$}}
\usq{ڈگری سے ریڈین میں تبدیل کریں: \lr{$315^{\circ}$}}
\usq{ڈگری سے ریڈین میں تبدیل کریں: \lr{$15^{\circ}15'$}}
\usq{ایک علاقے کی قوس کی لمبائی معلوم کریں اگر رداس \lr{10} سینٹی میٹر اور مرکزی زاویہ \lr{$\theta=60^{\circ}$} ہو۔}
\usq{ایک علاقے کا رقبہ معلوم کریں جس میں رداس \lr{r=8cm} اور مرکزی زاویہ \lr{$\theta=45^{\circ}$} ہے۔}
\usq{اگر کسی علاقے میں قوس کی لمبائی \lr{10} سینٹی میٹر، اور رداس \lr{5} سینٹی میٹر ہو تو قوس کے ذریعے بنائے گئے مرکزی زاویہ کی مقدار ریڈین اور ڈگری میں معلوم کریں۔}
\usq{اگر \lr{$\theta$} پہلے ربع میں ہے تو \lr{$\theta$} کی باقی تکونی نسبتیں معلوم کریں: \lr{$\sin\theta=\dfrac{2}{3}$}}
\usq{اگر \lr{$\theta$} پہلے ربع میں ہے تو \lr{$\theta$} کی باقی تکونی نسبتیں معلوم کریں: \lr{$\cos\theta=\dfrac{3}{4}$}}
\usq{اگر \lr{$\theta$} پہلے ربع میں ہے تو \lr{$\theta$} کی باقی تکونی نسبتیں معلوم کریں: \lr{$\tan\theta=\dfrac{1}{2}$}}
\usq{اگر \lr{$\theta$} پہلے ربع میں ہے تو \lr{$\theta$} کی باقی تکونی نسبتیں معلوم کریں: \lr{$\sec\theta=3$}}
\usq{اگر \lr{$\theta$} پہلے ربع میں ہے تو \lr{$\theta$} کی باقی تکونی نسبتیں معلوم کریں: \lr{$\cot\theta=\sqrt{\dfrac{3}{2}}$}}
\usq{درج ذیل تکونیاتی اکائی کو ثابت کریں: \lr{$(\sin\theta+\cos\theta)^{2}=1+2\sin\theta\cos\theta$}}
\usq{درج ذیل تکونیاتی اکائی کو ثابت کریں: \lr{$\dfrac{\cos\theta}{\sin\theta}=\dfrac{1}{\tan\theta}$}}
\usq{درج ذیل تکونیاتی اکائی کو ثابت کریں: \lr{$\cos^{2}\theta-\sin^{2}\theta=2\cos^{2}\theta-1$}}
\usq{درج ذیل تکونیاتی اکائی کو ثابت کریں: \lr{$\cos^{2}\theta-\sin^{2}\theta=1-2\sin^{2}\theta$}}
\usq{ثابت کریں: \lr{$(\sec^{2}\theta-1)\cos^{2}\theta=\sin^{2}\theta$}}
\usq{ثابت کریں: \lr{$\tan\theta+\cos\theta=\sec\theta\,\mathrm{cosec}\,\theta$}}
\usq{کیلکولیٹر کے استعمال کے بغیر درج ذیل تکونیاتی نسبت کی مقدار معلوم کریں: \lr{$\sin30^{\circ}$}}
\usq{کیلکولیٹر کے استعمال کے بغیر درج ذیل تکونیاتی نسبت کی مقدار معلوم کریں: \lr{$\cos30^{\circ}$}}
\usq{کیلکولیٹر کے استعمال کے بغیر درج ذیل تکونیاتی نسبت کی مقدار معلوم کریں: \lr{$\tan\dfrac{\pi}{6}$}}
\usq{کیلکولیٹر کے استعمال کے بغیر درج ذیل تکونیاتی نسبت کی مقدار معلوم کریں: \lr{$\tan60^{\circ}$}}
\usq{کیلکولیٹر کے استعمال کے بغیر درج ذیل تکونیاتی نسبت کی مقدار معلوم کریں: \lr{$\sec60^{\circ}$}}
\usq{کیلکولیٹر کے استعمال کے بغیر درج ذیل تکونیاتی نسبت کی مقدار معلوم کریں: \lr{$\cos\dfrac{\pi}{3}$}}
\usq{کیلکولیٹر کے استعمال کے بغیر درج ذیل تکونیاتی نسبت کی مقدار معلوم کریں: \lr{$\cot60^{\circ}$}}
\usq{کیلکولیٹر کے استعمال کے بغیر درج ذیل تکونیاتی نسبت کی مقدار معلوم کریں: \lr{$\sin60^{\circ}$}}
\usq{کیلکولیٹر کے استعمال کے بغیر درج ذیل تکونیاتی نسبت کی مقدار معلوم کریں: \lr{$\sec30^{\circ}$}}
\usq{کیلکولیٹر کے استعمال کے بغیر درج ذیل تکونیاتی نسبت کی مقدار معلوم کریں: \lr{$\mathrm{cosec}\,30^{\circ}$}}
\usq{کیلکولیٹر کے استعمال کے بغیر درج ذیل تکونیاتی نسبت کی مقدار معلوم کریں: \lr{$\sin45^{\circ}$}}
\usq{کیلکولیٹر کے استعمال کے بغیر درج ذیل تکونیاتی نسبت کی مقدار معلوم کریں: \lr{$\cos\dfrac{\pi}{4}$}}
\usq{پڑتال کریں: \lr{$2\sin60^{\circ}\cos60^{\circ}$}}
\usq{پڑتال کریں: \lr{$2\cos\dfrac{\pi}{3}\sin\dfrac{\pi}{3}$}}
\usq{پڑتال کریں: \lr{$2\sin45^{\circ}+2\cos45^{\circ}$}}
\usq{پڑتال کریں: \lr{$\cos60^{\circ}\cos30^{\circ}-\sin60^{\circ}\sin30^{\circ}$}}
\usq{پڑتال کریں: \lr{$\sin60^{\circ}\cos30^{\circ}-\cos60^{\circ}\sin30^{\circ}$}}
\usq{پڑتال کریں: \lr{$\cos60^{\circ}\cos30^{\circ}+\sin60^{\circ}\sin30^{\circ}$}}
\usq{پڑتال کریں: \lr{$\tan\dfrac{\pi}{6}\cot\dfrac{\pi}{6}+1$}}
\usq{درج ذیل باقاعدہ قائمۃ الزاویہ مثلث میں \lr{x,y,z} کی مقداریں معلوم کریں:}
\tfig{\coordinate(A)at(0,0);\coordinate(B)at(4,0);\coordinate(C)at(4,2.3);
\draw(A)--(B)--(C)--cycle;\draw(3.75,0)--(3.75,.25)--(4,.25);\draw(.9,0)arc(0:30:.9);\node at(1.25,.2){$30^\circ$};
\node[left]at(A){$A$};\node[below right]at(B){$B$};\node[above]at(C){$C$};
\node[below]at(2,0){$y=4$ cm};\node[right]at(4,1.1){$x$};\node[above left]at(2,1.15){$z$};}
\usq{درج ذیل باقاعدہ قائمۃ الزاویہ مثلث میں \lr{x,y,z} کی مقداریں معلوم کریں:}
\tfig{\coordinate(A)at(0,0);\coordinate(B)at(3,0);\coordinate(C)at(3,3);
\draw(A)--(B)--(C)--cycle;\draw(2.75,0)--(2.75,.25)--(3,.25);\draw(.9,0)arc(0:45:.9);\node at(1.15,.3){$\frac{\pi}{4}$};
\node[below left]at(A){$A$};\node[below right]at(B){$B$};\node[above]at(C){$C$};
\node[below]at(1.5,0){$x$};\node[above left]at(1.5,1.5){$z$};\node[rotate=90,anchor=south]at(3,1.5){$y=\sqrt{3}$ cm};}
\usq{درج ذیل باقاعدہ قائمۃ الزاویہ مثلث میں \lr{x,y,z} کی مقداریں معلوم کریں:}
\tfig{\coordinate(A)at(0,0);\coordinate(B)at(4,0);\coordinate(C)at(4,4);
\draw(A)--(B)--(C)--cycle;\draw(3.75,0)--(3.75,.25)--(4,.25);\draw(.9,0)arc(0:45:.9);\node at(1.3,.3){$45^\circ$};
\draw([shift={(4,4)}]-90:.6)arc(-90:-135:.6);
\node[left]at(A){$A$};\node[below right]at(B){$B$};\node[above]at(C){$C$};
\node[below]at(2,0){$y=4$ cm};\node[right]at(4,2){$x$};\node[above left]at(2,2){$z$};}
\usq{درج ذیل باقاعدہ قائمۃ الزاویہ مثلث میں \lr{x,y,z} کی مقداریں معلوم کریں:}
\tfig{\coordinate(C)at(0,0);\coordinate(B)at(2,0);\coordinate(A)at(2,3.46);
\draw(C)--(B)--(A)--cycle;\draw(1.75,0)--(1.75,.25)--(2,.25);\draw(.7,0)arc(0:60:.7);\node at(1.05,.35){$60^\circ$};
\node[left]at(C){$C$};\node[below right]at(B){$B$};\node[above]at(A){$A$};
\node[below]at(1,0){$x$};\node[right]at(2,1.7){$y$};\node[rotate=60,above]at(1,1.73){$z=2$ cm};}
\usq{درج ذیل مثلث کے نامعلوم اضلاع اور زاویے معلوم کریں۔}
\tfig{\coordinate(A)at(0,0);\coordinate(B)at(1.6,0);\coordinate(C)at(1.6,3);
\draw(A)--(B)--(C)--cycle;\draw(.5,0)arc(0:62:.5);\draw([shift={(C)}]-90:.5)arc(-90:-118:.5);
\node[left]at(A){$A$};\node[below right]at(B){$B$};\node[above]at(C){$C$};
\node[below]at(.8,0){$c=\sqrt{13}$ cm};\node[left,above]at(.8,1.5){$b$};\node[rotate=90,anchor=south]at(1.6,1.5){$a=\sqrt{3}$ cm};}
\usq{درج ذیل مثلث کے نامعلوم اضلاع اور زاویے معلوم کریں۔}
\tfig{\coordinate(A)at(0,0);\coordinate(B)at(4,0);\coordinate(C)at(4,3.3);
\draw(A)--(B)--(C)--cycle;\draw(3.75,0)--(3.75,.25)--(4,.25);
\node[left]at(A){$A$};\node[below right]at(B){$B$};\node[above]at(C){$C$};
\node[below]at(2,0){$c=4$ cm};\node[above left]at(2,1.65){$b$};\node[rotate=90,anchor=south]at(4,1.65){$a=4$ cm};}
\usq{ایک مربع کھیت کے ہر ضلع کی لمبائی \lr{60} میٹر ہے اس کھیت کے وتروں کی مقداریں معلوم کریں۔}
\usq{ہر ایک جزو میں \lr{x} کی قیمت معلوم کریں۔}
\tfig{\coordinate(B)at(0,2.5);\coordinate(C)at(4.5,2.5);\coordinate(A)at(0,0);
\draw(B)--(C)--(A)--cycle;\draw(0,2.25)--(.25,2.25)--(.25,2.5);
\node[above left]at(B){$B$};\node[above right]at(C){$C$};\node[below]at(A){$A$};
\node[above]at(2.25,2.5){12 cm};\node[left]at(0,1.25){$x$ cm};\node[rotate=-29,above]at(2.25,1.25){20 cm};}
\usq{ہر ایک جزو میں \lr{x} کی قیمت معلوم کریں۔}
\tfig{\coordinate(A)at(0,0);\coordinate(B)at(3,0);\coordinate(C)at(5,0);\coordinate(D)at(3,2.5);
\draw(A)--(C)--(D)--cycle;\draw(B)--(D);\draw(3,.25)--(3.25,.25)--(3.25,0);
\node[below left]at(A){$A$};\node[below]at(B){$B$};\node[below right]at(C){$C$};\node[above]at(D){$D$};
\node[above left,rotate=40]at(1.5,1.25){$x$ cm};\node[above right,rotate=-52]at(4,1.25){5 cm};\node[below]at(4,0){4 cm};
\node[below]at(3,-.45){$m\overline{AB}=m\overline{BC}$};}
\usq{مثلث \lr{ABC} حل کیجیے جس میں \lr{$m\angle A=60^{\circ}$}، \lr{$b=5$ cm} اور \lr{$m\angle B=90^{\circ}$}۔}
\usq{مثلث \lr{ABC} حل کریں۔ جب کہ \lr{$a=2$cm}، \lr{$b=2\sqrt{2}$ cm} اور \lr{$m\angle B=90^{\circ}$}}
\usq{زمین کی سطح پر \lr{40m} کے فاصلے سے کسی جھنڈے کی پوسٹ سے اوپر والی سطح کا زاویہ نزول \lr{$60^{\circ}$} ہے۔ پوسٹ کی بلندی معلوم کریں۔}
\usq{ایک درخت کی اونچائی \lr{72m} ہے۔ اس درخت سے \lr{10m} کے فاصلے سے اس کے اوپر والے سرے کا زاویہ نزول معلوم کریں۔}
\usq{ایک سیڑھی زمین کے ساتھ \lr{$60^{\circ}$} کا زاویہ بناتی ہے اور دیوار کے ساتھ \lr{10m} کی بلندی تک پہنچتی ہے۔ اس سیڑھی کی لمبائی معلوم کریں۔}
\usq{اگر \lr{300m} اونچے مینار کا سایہ \lr{450m} ہو تو سورج کے زاویہ نزول کی مقدار معلوم کریں۔}
\usq{دیے گئے زاویوں کو ریڈین میں تبدیل کریں اور جواب \lr{$\pi$} کی شکل میں لکھیں: \lr{$255^{\circ}$}}
\usq{دیے گئے زاویوں کو ریڈین میں تبدیل کریں اور جواب \lr{$\pi$} کی شکل میں لکھیں: \lr{$75^{\circ}45'$}}
\usq{دیے گئے زاویوں کو ریڈین میں تبدیل کریں اور جواب \lr{$\pi$} کی شکل میں لکھیں: \lr{$142.5^{\circ}$}}
\usq{ریڈین سے ڈگری میں تبدیل کریں اور جواب ڈگری، منٹ اور سیکنڈ میں لکھیں: \lr{$\dfrac{17\pi}{24}$}}
\usq{ریڈین سے ڈگری میں تبدیل کریں اور جواب ڈگری، منٹ اور سیکنڈ میں لکھیں: \lr{$\dfrac{11\pi}{16}$}}
\usq{درج ذیل تکونیاتی اکائی کو ثابت کریں: \lr{$\dfrac{\sin\theta}{1-\cos\theta}=\dfrac{1+\cos\theta}{\sin\theta}$}}
\usq{نقطہ \lr{$P(x,y)$} کے مستوی میں مقام کی وضاحت کریں، جس کے لیے: \lr{$x>0$}}
\usq{نقطہ \lr{$P(x,y)$} کے مستوی میں مقام کی وضاحت کریں، جس کے لیے: \lr{$x>0$} اور \lr{$y>0$}}
\usq{نقطہ \lr{$P(x,y)$} کے مستوی میں مقام کی وضاحت کریں، جس کے لیے: \lr{$x=0$}}
\usq{نقطہ \lr{$P(x,y)$} کے مستوی میں مقام کی وضاحت کریں، جس کے لیے: \lr{$y=0$}}
\usq{نقطہ \lr{$P(x,y)$} کے مستوی میں مقام کی وضاحت کریں، جس کے لیے: \lr{$x\geq0$} اور \lr{$y<0$}}
\usq{نقطہ \lr{$P(x,y)$} کے مستوی میں مقام کی وضاحت کریں، جس کے لیے: \lr{$x=0,\ y=0$}}
\usq{نقطہ \lr{$P(x,y)$} کے مستوی میں مقام کی وضاحت کریں، جس کے لیے: \lr{$|x|=|y|$}}
\usq{نقطہ \lr{$P(x,y)$} کے مستوی میں مقام کی وضاحت کریں، جس کے لیے: \lr{$x\geq3$}}
\usq{نقطہ \lr{$P(x,y)$} کے مستوی میں مقام کی وضاحت کریں، جس کے لیے: \lr{$y>2$}}
\usq{نقطہ \lr{$P(x,y)$} کے مستوی میں مقام کی وضاحت کریں، جس کے لیے: \lr{x} اور \lr{y} کی مخالف علامتیں ہیں۔}
\usq{درج ذیل نقاط کے درمیان فاصلہ معلوم کریں: \lr{$A(6,7),\ B(0,-2)$}}
\usq{درج ذیل نقاط کے درمیان فاصلہ معلوم کریں: \lr{$C(-5,-2),\ D(3,2)$}}
\usq{درج ذیل نقاط کے درمیان فاصلہ معلوم کریں: \lr{$L(0,3),\ M(-2,-4)$}}
\usq{درج ذیل نقاط کے درمیان فاصلہ معلوم کریں: \lr{$P(-8,-7),\ Q(0,0)$}}
\usq{درج ذیل معلوم کریں: دیے گئے دو نقاط کے درمیان فاصلہ۔}
\usq{درج ذیل میں معلوم کریں: دو نقاط کو ملانے والے خط کا وسطی نقطہ: \lr{$A(3,1),\ B(-2,-4)$}}
\usq{درج ذیل میں معلوم کریں: دو نقاط کو ملانے والے خط کا وسطی نقطہ: \lr{$A(-8,3),\ (2,-1)$}}
\usq{درج ذیل میں معلوم کریں: دو نقاط کو ملانے والے خط کا وسطی نقطہ: \lr{$A\!\left(-\sqrt{5},-\dfrac{1}{3}\right),\ B(-3\sqrt{5},\ 5)$}}
\usq{کیا دیا گیا نقطہ مبدا سے \lr{15} اکائیوں کے فاصلہ پر ہے؟ \lr{$(\sqrt{176},\ 7)$}}
\usq{کیا دیا گیا نقطہ مبدا سے \lr{15} اکائیوں کے فاصلہ پر ہے؟ \lr{$(10,-10)$}}
\usq{کیا دیا گیا نقطہ مبدا سے \lr{15} اکائیوں کے فاصلہ پر ہے؟ \lr{$(1,15)$}}
\usq{جیومیٹری کی تاریخ پر مختصر لکھیں۔}
\usq{مبدا کی تعریف کریں۔}
\usq{محدداتی محور کی تعریف کریں۔}
\usq{ربع کی تعریف کریں۔}
\usq{دو رخی محدداتی نظام کی تعریف کریں۔}
\usq{فاصلہ کا فارمولا لکھیں۔}
\usq{درمیانی نقطہ کا فارمولا لکھیں۔}
\usq{دیے گئے نقاط کے درمیان فاصلہ معلوم کریں: \lr{$A(5,6),\ B(5,-2)$}}
\usq{ثابت کریں کہ نقاط \lr{$A(-1,2),\ B(7,5)$} اور \lr{$C(2,-6)$} قائمۃ الزاویہ مثلث کے راس ہیں۔}
\usq{نقطہ \lr{$C(-5,3)$} دائرہ کا مرکز اور \lr{$P(7,-2)$} دائرہ پر واقع ہے۔ دائرہ کا رداس معلوم کریں۔}
\usq{\lr{$A(2,3)$} اور \lr{$B(8,7)$} کو ملانے والے خط کا وسطی نقطہ معلوم کریں۔}
\usq{دیے گئے نقاط سے بننے والے خط کی ڈھلوان اور جھکاؤ معلوم کریں: \lr{$(-2,4);\ (5,11)$}}
\usq{دیے گئے نقاط سے بننے والے خط کی ڈھلوان اور جھکاؤ معلوم کریں: \lr{$(3,-2);\ (2,7)$}}
\usq{دیے گئے نقاط سے بننے والے خط کی ڈھلوان اور جھکاؤ معلوم کریں: \lr{$(4,6);\ (4,8)$}}
\usq{ڈھلوان کے ذریعے یہ ظاہر کریں کہ درج ذیل نقاط ایک ہی خط پر ہیں: \lr{$A(-1,-3);\ B(1,5);\ C(2,9)$}}
\usq{ڈھلوان کے ذریعے یہ ظاہر کریں کہ درج ذیل نقاط ایک ہی خط پر ہیں: \lr{$P(4,-5);\ Q(7,5);\ R(10,15)$}}
\usq{ڈھلوان کے ذریعے یہ ظاہر کریں کہ درج ذیل نقاط ایک ہی خط پر ہیں: \lr{$L(-4,6);\ M(3,8);\ N(10,10)$}}
\usq{ڈھلوان کے ذریعے یہ ظاہر کریں کہ درج ذیل نقاط ایک ہی خط پر ہیں: \lr{$X(a,2b);\ Y(c,a+b);\ Z(2c-a,2a)$}}
\usq{\lr{k} معلوم کریں جب کہ \lr{$A(7,3);\ B(k,-6)$} کو ملانے والا خط اور \lr{$C(-4,5);\ D(-6,4)$} کو ملانے والا خط متوازی ہوں}
\usq{\lr{k} معلوم کریں جب کہ \lr{$A(7,3);\ B(k,-6)$} کو ملانے والا خط اور \lr{$C(-4,5);\ D(-6,4)$} کو ملانے والا خط عمودی ہوں۔}
\usq{مساوات معلوم کریں: \lr{$(7,-9)$} میں سے گزرنے والا افقی خط}
\usq{مساوات معلوم کریں: \lr{$(-5,3)$} میں سے گزرنے والا عمودی خط ہو}
\usq{مساوات معلوم کریں: \lr{$A(-6,5)$} سے گزرنے والا خط جس کی ڈھلوان \lr{7} ہے۔}
\usq{مساوات معلوم کریں: \lr{$(8,-3)$} سے گزرنے والا خط جس کی ڈھلوان \lr{0} ہے۔}
\usq{مساوات معلوم کریں: \lr{$(-8,5)$} سے گزرنے والا خط جس کی ڈھلوان غیر واضح ہو۔}
\usq{مساوات معلوم کریں: \lr{$(-5,-3)$} اور \lr{$(9,-1)$} سے گزرنے والا خط۔}
\usq{مساوات معلوم کریں: \lr{y}- قاطع: \lr{$-7$} اور ڈھلوان \lr{5}}
\usq{مساوات معلوم کریں: \lr{x} قاطع: \lr{$-3$} اور \lr{y} قاطع \lr{4}}
\usq{مساوات معلوم کریں: \lr{x} قاطع: \lr{$-9$} اور ڈھلوان \lr{$-4$}}
\usq{ڈھلوان کی تعریف کریں۔}
\usq{\lr{x}-قطع کی تعریف کریں۔}
\usq{\lr{y}-قطع کی تعریف کریں۔}
\usq{ثابت کریں کہ نقاط \lr{$A(-3,6),\ B(3,2)$} اور \lr{$C(6,0)$} ہم خط ہیں۔}
\usq{ثابت کریں کہ نقاط \lr{$A(1,1),\ B(4,5)$} اور \lr{$C(12,-1)$} ایک قائمۃ الزاویہ مثلث کے راس ہیں۔}
\usq{ایسے خط کی مساوات معلوم کریں جو \lr{$(-2,1)$} اور \lr{$(6,-4)$} میں سے گزرتا ہو۔}
\usq{خط کی ایک مساوات لکھیں جو \lr{x}-محور کو \lr{$(2,0)$} اور \lr{y}-محور کو \lr{$(0,-4)$} پر قطع کرتی ہے۔}
\usq{نقطہ \lr{$P(2,3)$} میں سے گزرنے والے خط کی مساوات معلوم کریں جو پہلے ربع میں محدداتی محور کے ساتھ ایک متساوی الساقین مثلث بناتی ہو۔}
\usq{\lr{$5x-12y+39=0$} کو درج ذیل شکل میں تبدیل کریں: ڈھلوان قاطع شکل}
\usq{\lr{$5x-12y+39=0$} کو درج ذیل شکل میں تبدیل کریں: دو قاطع شکل}
\usq{\lr{$5x-12y+39=0$} کو درج ذیل شکل میں تبدیل کریں: عام شکل}
\usq{\lr{$5x-12y+39=0$} کو درج ذیل شکل میں تبدیل کریں: نقطہ ڈھلوان شکل}
\usq{\lr{$5x-12y+39=0$} کو درج ذیل شکل میں تبدیل کریں: دو نقطی شکل}
\usq{\lr{$5x-12y+39=0$} کو درج ذیل شکل میں تبدیل کریں: سمیٹرک شکل}
\usq{دو دوستوں کے گھروں کو ایک گرڈ پر محدادت \lr{$(2,6)$} اور \lr{$(9,12)$} کے ذریعے ظاہر کیا گیا ہے۔ اگر گرڈ کی اکائیاں کلومیٹر کو ظاہر کرتی ہوں تو ان کے گھروں کے درمیان فاصلہ معلوم کریں۔}
\usq{ایک سیدھی پگڈنڈی پر غور کریں جو نقطہ \lr{$(5,7)$} سے شروع ہوتی ہے اور نقطہ \lr{$(15,3)$} پر ختم ہوتی ہے۔ اس کے وسطی نقطہ کے محددات معلوم کریں۔}
\usq{ریس ٹریک کے آغاز اور اختتامی نقاط کے محددات \lr{$(3,9)$} اور \lr{$(9,13)$} ہیں۔ ٹریک کا وسطی نقطہ کیا ہے؟}
\usq{سڑک پر دو نقاط \lr{$A(3,4)$} اور \lr{$B(7,10)$} ہیں۔ سڑک کے وسطی نقطہ کے محددات معلوم کریں۔}
\usq{ایوی ایشن کیا ہے؟}
\usq{نیوی گیشن کیا ہے؟}
\usq{عرض بلد کیا ہے؟}
\usq{طول بلد کیا ہے؟}
\usq{نقشے پر ٹاؤن \lr{A} کے محددات \lr{$(2,3)$} اور ٹاؤن \lr{B} کے محددات \lr{$(-4,-1)$} ہیں۔ دونوں ٹاؤن کے درمیان فاصلہ معلوم کریں}
\usq{فرض کریں دو شہروں شہر \lr{A} اور شہر \lr{B} کے محددات \lr{$(3,4)$} اور \lr{$(7,1)$} نقشے پر ظاہر کیے گئے ہیں۔ دونوں شہروں کے درمیان فاصلہ معلوم کریں۔}
\usq{ایک انجینئر دریا کے کنارے پر دو مقامات کے درمیان پل بنا رہا ہے۔ فرض کریں کہ \lr{$(2,5)$} اور \lr{$(8,9)$} دو مقامات کے نقاط ہیں جہاں سے پل شروع اور ختم ہوگا۔ وسطی نقطہ کے محددات معلوم کریں جو پل کے مرکز کو ظاہر کرتے ہیں۔}
\usq{ایک پائلٹ کو شہر \lr{A} \lr{$(50,60)$} سے شہر \lr{B} \lr{$(120,150)$} تک سفر کرنے کی ضرورت ہے۔ ہیڈنگ زاویہ کا تعین کریں جو ہوائی جہاز کو مشرق کی سمت لے جانا چاہیے۔}
\usq{عبدالہادی نقطہ \lr{$A$(Latitude $10^{\circ}$ N, Longitude $50^{\circ}$ E)} سے نقطہ \lr{$B$(Latitude $20^{\circ}$ N, Longitude $60^{\circ}$ E)} تک سفر کر رہا ہے۔ عرض بلد اور طول بلد کے لحاظ سے اس کے سفر کا وسطی نقطہ معلوم کریں۔}
\usq{ایک زمین کی تزئین و آرائش کرنے والا \lr{$P(2,3)$} سے \lr{$Q(8,9)$} تک سیدھا راستہ ڈائزن کر رہا ہے۔ راستے کی لمبائی معلوم کریں}
\usq{محدداتی مستوی میں نقاط \lr{$A(2,3)$} اور \lr{$B(7,8)$} کے درمیان فاصلہ معلوم کریں۔}
\usq{نقاط \lr{$(4,-2)$} اور \lr{$(-6,3)$} سے مل کر بننے والے قطعہ خط کا وسطی نقطہ معلوم کریں۔}
\usq{نقاط \lr{$(1,2)$} اور \lr{$(4,6)$} سے گزرنے والے خط کی ڈھلوان معلوم کریں۔}
\usq{اگر دو خطوط متوازی ہوں اور ایک خط کی ڈھلوان \lr{$\dfrac{2}{3}$} ہو تو دوسرے خط کی ڈھلوان معلوم کریں۔}
\usq{ایک ہوائی جہاز کو شہر \lr{A} پر محددات \lr{$(12,5)$} سے شہر \lr{B} پر محددات \lr{$(8,-4)$} تک پرواز کرنے کی ضرورت ہے۔ ان دو شہروں کے درمیان فاصلہ معلوم کریں۔}
\usq{زمین کی تزئین و آرائش کے منصوبے میں راستے کے محددات \lr{$(2,3)$} اور \lr{$(10,7)$} ہیں۔ اس راستے کے وسطی نقاط کے محددات معلوم کریں۔}
\usq{\lr{$-3$} کے میلان اور \lr{2} کے \lr{y}- قاطع والے خط کے لیے مساوات لکھیں۔ ڈھلوان قاطع شکل}
\usq{\lr{$-3$} کے میلان اور \lr{2} کے \lr{y}- قاطع والے خط کے لیے مساوات لکھیں۔ نقطہ \lr{$(1,2)$} کا استعمال کرتے ہوئے نقطہ ڈھلوان شکل}
\usq{\lr{$-3$} کے میلان اور \lr{2} کے \lr{y}- قاطع والے خط کے لیے مساوات لکھیں۔ نقاط \lr{$(1,2)$} اور \lr{$(4,-7)$} کا استعمال کرتے ہوئے دو نقطی شکل}
\usq{\lr{$-3$} کے میلان اور \lr{2} کے \lr{y}- قاطع والے خط کے لیے مساوات لکھیں۔ قاطع شکل}
\usq{\lr{$-3$} کے میلان اور \lr{2} کے \lr{y}- قاطع والے خط کے لیے مساوات لکھیں۔ سمیٹرک شکل}
\usq{\lr{$-3$} کے میلان اور \lr{2} کے \lr{y}- قاطع والے خط کے لیے مساوات لکھیں۔ عام شکل}
\usq{درج ذیل مشروط کے الٹ، معکوس اور ضد مثبت لکھیے۔ \lr{$\sim p\rightarrow q$}}
\usq{درج ذیل مشروط کے الٹ، معکوس اور ضد مثبت لکھیے۔ \lr{$q\rightarrow p$}}
\usq{درج ذیل مشروط کے الٹ، معکوس اور ضد مثبت لکھیے۔ \lr{$\sim q\rightarrow\sim p$}}
\usq{درج ذیل مشروط کے الٹ، معکوس اور ضد مثبت لکھیے۔ \lr{$\sim q\rightarrow\sim p$}}
\usq{درج ذیل کا صدق جدول لکھیں: \lr{$\sim(p\vee q)\vee(\sim q)$}}
\usq{درج ذیل کا صدق جدول لکھیں: \lr{$\sim(\sim q\vee\sim p)$}}
\usq{درج ذیل کا صدق جدول لکھیں: \lr{$(p\ \vee\ q)\leftrightarrow(p\ \wedge\ p)$}}
\usq{ایک ریاضیاتی بیان اور اس کے ثبوت میں فرق واضح کریں۔ دو مثالیں دیجیے۔}
\usq{ایک مسئلہ اور اصول متعارفہ میں کیا فرق ہے؟ ہر ایک کہ مثال دیجیے۔}
\usq{ریاضیاتی ثبوتوں میں منطقی استدلال کی اہمیت کیا ہے؟ اپنی بات کو واضح کرنے کے لیے ایک مثال دیجیے۔}
\usq{یہ بتائیں کہ دیا گیا بیان ایک اصول متعارفہ، قیاس یا مسئلہ ہے اور اپنی دلیل کی وضاحت کریں۔ کوئی دو نقاط میں سے صرف ایک خط گزر تا ہے۔}
\usq{یہ بتائیں کہ دیا گیا بیان ایک اصول متعارفہ، قیاس یا مسئلہ ہے اور اپنی دلیل کی وضاحت کریں۔ ہر جعت عدد \lr{2} سے بڑے دو مفرد اعداد کے مجموعے کے طور پر لکھا جاسکتا ہے۔}
\usq{یہ بتائیں کہ دیا گیا بیان ایک اصول متعارفہ، قیاس یا مسئلہ ہے اور اپنی دلیل کی وضاحت کریں۔ مثلث کے زاویوں کا مجموعہ \lr{$180^{\circ}$} ہوتا ہے۔}
\usq{فرض کریں \lr{x} ایک صحیح عدد ہے۔ اگر \lr{x} ایک طاق عدد ہو تو \lr{$9x+4$} ایک طاق عدد ہوگا۔}
\usq{فرض کریں \lr{x} ایک صحیح عدد ہے۔ اگر \lr{x} ایک طاق عدد ہو تو \lr{$7x+5$} ایک جفت عدد ہوگا۔}
\usq{اگر \lr{x} اور \lr{y} مثبت حقیقی اعداد ہوں اور \lr{$x^{2}<y^{2}$} ہو تو \lr{$x<y$}۔}
\usq{ثابت کریں کہ مثلث کے اندرونی زاویوں کا مجموعہ \lr{$180^{\circ}$} ہوتا ہے۔}
\usq{\lr{If a,b and c are non-zero real numbers, prove that: $\dfrac{a}{b}\ \dfrac{c}{d}\Leftrightarrow ad=bc$}}
\usq{اگر \lr{a,b} اور \lr{c} غیر صفر حقیقی اعداد ہوں تو ثابت کریں کہ: \lr{$\dfrac{a}{b}\cdot\dfrac{c}{d}=\dfrac{ac}{bd}$}}
\usq{اگر \lr{a,b} اور \lr{c} غیر صفر حقیقی اعداد ہوں تو ثابت کریں کہ: \lr{$\dfrac{a}{b}+\dfrac{c}{b}=\dfrac{a+c}{b}$}}
\usq{استقراء کی تعریف کریں۔}
\usq{استخراج کی تعریف کریں۔}
\usq{بیان کی تعریف کریں۔}
\usq{نفی کیا ہے؟}
\usq{کنجکشن کیا ہے؟}
\usq{ڈس جنکشن کیا ہے؟}
\usq{برابری کی تعریف کریں۔}
\usq{مسئلہ کی تعریف کریں۔}
\usq{قیاس کیا ہے؟}
\usq{اصول متعارفہ کی تعریف کریں۔}
\usq{اقلیدس کا اصول متعارفہ کیا ہے؟}
\usq{پیانو کا اصول متعارفہ کیا ہے؟}
\usq{استخراجی ثبوت کیا ہے؟}
\usq{آیا کہ درج ذیل بیان درست ہے یا غلط: لاہور پنجاب کا دارالحکومت ہے اور کوئٹہ بلوچستان کا دارالحکومت ہے۔}
\usq{آیا کہ درج ذیل بیان درست ہے یا غلط: \lr{$4<5\wedge8<10$}}
\usq{آیا کہ درج ذیل بیان درست ہے یا غلط: \lr{$2+2=3\wedge6+6=10$}}
\usq{\lr{10} ایک مثبت عدد ہے یا \lr{0} ایک ناطق عدد ہے۔ اس ڈس جنکشن کی درست قیمت معلوم کریں۔}
\usq{مثلث کے دو قائمۃ الزاویے ہوسکتے ہیں یالاہور سندھ کا دارالحکومت ہے۔ اس ڈس جنکشن کی درست قیمت معلوم کریں۔}
\usq{ثابت کریں کہ کسی بھی کائناتی سیٹ میں خالی سیٹ \lr{$\phi$} کسی بھی سیٹ \lr{A} کا تختی سیٹ ہوتا ہے۔}
\usq{درج ذیل ریاضیاتی بیان کو ثابت کریں: اگر \lr{x} ایک طاق عدد ہو تو \lr{$x^{2}$} بھی ایک طاق عدد ہوتا ہے۔}
\usq{درج ذیل ریاضیاتی بیان کو ثابت کریں: دو طاق اعداد کا مجموعہ ایک جفت عدد ہوتا ہے۔}
\usq{معلوم کریں کہ کیا دیے گئے اجسام متشابہ ہیں؟ تمام لمبائیاں \lr{cm} میں ہیں؟}
\pfig{fig665}{5.5cm}
\usq{مثلث \lr{ABC} میں اضلاع کی لمبائیاں \lr{$m\overline{AB}=6$ cm}، \lr{$m\overline{BC}=9$ cm} اور \lr{$m\overline{CA}=12$ cm} اور مثلث \lr{DEF} میں اضلاع کی لمبائیاں \lr{$m\overline{DE}=10.5$ cm}، \lr{$m\overline{EF}=15.75$ cm} اور \lr{$m\overline{FD}=21$ cm} ہیں۔ ثابت کریں دونوں مثلثیں متشابہ ہیں۔}
\usq{دی گئی شکل میں \lr{x} کی قیمت معلوم کریں۔}
\pfig{fig667}{5cm}
\usq{دی گئی شکل میں \lr{x} کی قیمت معلوم کریں۔}
\pfig{fig668}{5cm}
\usq{دی گئی شکل میں \lr{x} کی قیمت معلوم کریں۔}
\pfig{fig669}{4.5cm}
\usq{دی گئی شکل میں \lr{x,y} اور \lr{z} کی قیمتیں معلوم کریں۔}
\pfig{fig670}{4.5cm}
\usq{اگر متناظرہ دو اضلاع کا جوڑا دوسرے کے متوازی ہو تو دو مثلثیں متشابہ ہوں گی جیسا کہ دکھایا گیا ہے۔ مثلاً دی ہوئی شکل میں \lr{$\overline{AB}$}، \lr{$\overline{CD}$} کے متوازی ہے۔}
\usq{منظم مثمن کا احاطہ \lr{48cm} ہے۔ ایک دوسری مثمن کے اضلاع کی لمبائیاں پہلی مثمن کا \lr{1.2} گنا ہیں۔ دوسری منظم مثمن کے اضلاع کی لمبائیاں معلوم کریں۔}
\usq{متشابہ اشکال کے رقبوں کی نسبت معلوم کریں اگر ان کی متناظرہ لمبائیوں کی نسبت درج ذیل ہوں: \lr{1:3}}
\usq{متشابہ اشکال کے رقبوں کی نسبت معلوم کریں اگر ان کی متناظرہ لمبائیوں کی نسبت درج ذیل ہوں: \lr{3:4}}
\usq{متشابہ اشکال کے رقبوں کی نسبت معلوم کریں اگر ان کی متناظرہ لمبائیوں کی نسبت درج ذیل ہوں: \lr{2:7}}
\usq{متشابہ اشکال کے رقبوں کی نسبت معلوم کریں اگر ان کی متناظرہ لمبائیوں کی نسبت درج ذیل ہوں: \lr{8:9}}
\usq{دو کروں کے رداسوں میں نسبت \lr{3:4} ہے۔ ان کے حجموں میں نسبت معلوم کریں۔}
\usq{دو منظم چار اضلاع کے حجموں میں نسبت \lr{8:27} ہے۔ ان کے اضلاع کی لمبائیوں میں نسبت معلوم کریں۔}
\usq{دو مخروطوں کے حجموں میں نسبت \lr{64:125} ہے۔ نسبت کیا ہوگی؟ ان کی اونچائی میں}
\usq{دو مخروطوں کے حجموں میں نسبت \lr{64:125} ہے۔ نسبت کیا ہوگی؟ قاعدے کے رقبوں میں}
\usq{دیے گئے متشابہ اجسام میں نامعلوم مقداریں معلوم کریں۔}
\pfig{fig681}{6.5cm}
\usq{دیے گئے متشابہ اجسام میں نامعلوم مقداریں معلوم کریں۔}
\pfig{fig682}{5.5cm}
\usq{دیے گئے متشابہ اجسام میں نامعلوم مقداریں معلوم کریں۔}
\pfig{fig683}{6.5cm}
\usq{دیے گئے متشابہ اجسام میں نامعلوم مقداریں معلوم کریں۔}
\pfig{fig684}{6.5cm}
\usq{ڈیکاگان کے اندرونی زاویوں کا مجموعہ کتنا ہوگا؟}
\usq{منظم مخمس کے ہر بیرونی زوایے کی مقدار کتنی ہے؟}
\usq{منظم مسدس کے ہر اندرونی زاویے کی مقدار معلوم کریں۔}
\usq{اگر کسی کثیر الاضلاع کے اندرونی زاویوں کا مجموعہ \lr{$1260^{\circ}$} ہو تو اس کثیر اضلاع کے کتنے اضلاع ہوں گے؟}
\usq{ایک متوازی الاضلاع \lr{ABCD} میں اگر \lr{$m\angle BAD=45^{\circ}$} اور \lr{$m\overline{AB}=10$ cm}، \lr{$m\overline{AD}=6$cm} ہو تو اس کا رقبہ معلوم کریں۔}
\usq{ایک مستطیلی دیوار کی لمبائی \lr{10m} اور چوڑائی \lr{120m} ہے۔ اس دیوار کو رنگ کرنے کے لیے کتنے گیلن کافی ہیں؟ اگر ایک گیلن \lr{$35$ m$^{2}$} کے لیے کافی ہو۔}
\usq{ایک مستطیلی فرش میٹر \lr{12} بائے \lr{15} میٹر ہے۔ اس فرش کو مکمل کرنے کے لیے \lr{1} بائے \lr{1} میٹر کی کتنی ٹائلیں درکار ہوں گی؟}
\usq{ایک مستطیلی دیوار کی اونچائی \lr{10} میٹر اور چوڑائی \lr{4} میٹر ہے۔ اگر \lr{1} لیٹر رنگ \lr{$7$ m$^{2}$} کے لیے کافی ہو تو اس دیوار کے لیے کتنے لیٹر رنگ درکار ہوگا؟}
\usq{ایک \lr{s} ضلع والی لمبائی مساوی الاضلاع مثلث \lr{ABC} کا رقبہ معلوم کریں۔}
\usq{درج ذیل یک درجی تفاعل کا گراف بنائیں: \lr{$y=3x-5$}}
\usq{درج ذیل یک درجی تفاعل کا گراف بنائیں: \lr{$y=-2x+8$}}
\usq{درج ذیل یک درجی تفاعل کا گراف بنائیں: \lr{$y=0.5x-1$}}
\usq{درج ذیل تفاعل کا گراف بنائیں: \lr{$y=4^{x}$}}
\usq{درج ذیل تفاعل کا گراف بنائیں: \lr{$y=5^{-x}$}}
\usq{درج ذیل تفاعل کا گراف بنائیں: \lr{$y=\dfrac{1}{x-3},\ x\neq3$}}
\usq{درج ذیل تفاعل کا گراف بنائیں: \lr{$y=\dfrac{2}{x}+3,\ x\neq0$}}
\usq{درج ذیل تفاعل کا گراف بنائیں: \lr{$y=x^{\frac{1}{2}}$}}
\usq{درج ذیل تفاعل کا گراف بنائیں: \lr{$y=3x^{\frac{1}{3}}$}}
\usq{درج ذیل تفاعل کا گراف بنائیں: \lr{$y=2x^{-2}$}}
\usq{گراف کیا ہوتے ہیں؟}
\usq{دو درجی تفاعل کی تعریف کریں۔}
\usq{سہ درجی تفاعل کی تعریف کریں۔}
\usq{معکوس تفاعل کی تعریف کریں۔}
\usq{\lr{$y=2x-1$} کا گراف بنائیں۔}
\usq{ایک ہی شکل میں \lr{$y=x^{2}$} اور \lr{$y=-x^{2}$} کے گراف بنائیں۔}
\usq{\lr{$-3\leq x\leq6$} کے لیے \lr{$y=2x^{2}-7x-9$} کا گراف بنائیں۔}
\usq{\lr{$-3\leq x\leq3$} کے لیے درج ذیل سہ رتبی تفاعل کا گراف بنائیں: \lr{$y=-x^{3}$}}
\usq{\lr{$-2\leq x\leq3$} کے لیے \lr{$y=2x^{3}-3x^{2}+x-5$} کا گراف بنائیں۔}
\usq{درج ذیل معکوس تفاعل کا گراف بنائیں: \lr{$y=\dfrac{1}{x-0.5},\ x\neq0.5$}}
\usq{\lr{$-6\leq x\leq6$} کے لیے قوت نمائی تفاعل \lr{$y=2^{x}$} کا گراف بنائیں۔}
\usq{قوت نمائی تفاعل \lr{$y=e^{x}$} کا گراف بنائیں۔}
\usq{\lr{$-3\leq x\leq3$} کے لیے \lr{$y=x^{3}$} کا گراف بنائیں۔}
\usq{\lr{$y=x^{-1}$} کا گراف بنائیں۔}
\usq{\lr{$y=2x^{\frac{1}{5}}$} کا گراف بنائیں۔}
\usq{ایک کمپنی سکول بیگ تیار کرتی ہے۔ \lr{x} بیگ تیار کرنے کی لاگت \lr{C(x)=1200 + 800x} ہے اور \lr{x} بیگ کی فروخت سے حاصل ہونے والی آمدنی \lr{R(x)=50x} ہے۔ مساوی نقطہ معلوم کریں۔}
\usq{کسی شے کی طلب اور رسد کے تفاعل مساوات کے ذریعے دیے گئے ہیں۔ \lr{$P_{d}=400-5Q,\ P_{s}=3Q+24$} وقفہ \lr{Q=0} سے \lr{Q=300} پر ہر تفاعل کا گراف بنائیں۔}
\usq{ایک کمپنی سکول بیگ تیار کرتی ہے۔ \lr{x} بیگ تیار کرنے کی لاگت \lr{C(x)=1200 + 800x} ہے اور \lr{x} بیگ کی فروخت سے حاصل ہونے والی آمدنی \lr{R(x)=50x} ہے۔ \lr{250} بیگ فروخت ہونے پر نفع اور نقصان کا تعین کریں۔}
\usq{شاہد کی تنخواہ \lr{S(x)} روپوں میں درج ذیل کلیے پر مبنی ہے: \lr{S(x)=45000 + 4500x} جہاں \lr{x} ان سالوں کی تعداد ہے جو اس نے کمپنی میں کام کیا۔ \lr{$0\leq x\leq5$} کے لیے تنخواہ کے تفاعل کے گراف بنائیں اور اس کی تشریح کریں۔}
\usq{ایک کمپنی سکول بیگ تیار کرتی ہے۔ \lr{x} بیگ تیار کرنے کی لاگت \lr{C(x)=1200 + 800x} ہے اور \lr{x} بیگ کی فروخت سے حاصل ہونے والی آمدنی \lr{R(x)=50x} ہے۔ دونوں تفاعل کے گراف بنائیں اور مساوی نقطہ کی نشاندہی کریں۔}
\usq{درج ذیل تفاعل کا گراف بنائیں: \lr{$y=3^{-x}$}، \lr{x} کی قدروں \lr{$-2$} سے \lr{4} تک}
\usq{درج ذیل کا گراف بنائیں: \lr{$y=\dfrac{2}{x+7},\ x\neq-7$}}
\usq{\lr{$-5$} سے \lr{5} تک \lr{x} کے لیے درج ذیل کا گراف بنائیں: \lr{$y=x^{2}-3$}}
\usq{\lr{$-5$} سے \lr{5} تک \lr{x} کے لیے درج ذیل کا گراف بنائیں: \lr{$y=15-x^{2}$}}
\usq{لوکس سے کیا مراد ہے؟}
\usq{مختلف الاضلاع مثلث کیا ہے؟}
\usq{مساوی الساقین مثلث کیا ہے؟}
\usq{امساوی الاضلاع مثلث کیا ہے؟}
\usq{حادہ زاویہ مثلث کیا ہے؟}
\usq{منفرجہ زاویہ مثلث کیا ہے؟}
\usq{قائمہ زاویہ مثلث کیا ہے؟}
\usq{ض۔ض۔ز کیا ہے؟}
\usq{عمودی ناصف کی تعریف کریں۔}
\usq{ہم نقطہ کی تعریف کریں۔}
\usq{محاصرہ مرکز کی تعریف کریں۔}
\usq{مرکز نما کی تعریف کریں۔}
\usq{مثلث کے زاویہ کا ناصف کیا ہے؟}
\usq{محصور مرکز کی تعریف کریں۔}
\usq{ارتفاع کی تعریف کریں۔}
\usq{عمودی مرکز کی تعریف کریں۔}
\usq{وسطانیہ کی تعریف کریں۔}
\usq{لوکائی کیا ہے؟}
\usq{دائرہ کی تعریف کریں۔}
\usq{زاویہ کا ناصف کی تعریف کریں۔}
\usqd{کیمسٹری کے طلبہ کے ایک گروپ نے کیمیائی لیبارٹری میں درج ذیل نمبر حاصل کیے:}{آخری جماعت کی بالائی حد کیا ہے؟}
\usqd{کیمسٹری کے طلبہ کے ایک گروپ نے کیمیائی لیبارٹری میں درج ذیل نمبر حاصل کیے:}{جماعتی حد \lr{(39-43)} میں زیریں جماعتی حد کیا ہے؟}
\usqd{کیمسٹری کے طلبہ کے ایک گروپ نے کیمیائی لیبارٹری میں درج ذیل نمبر حاصل کیے:}{جماعتی حد \lr{(34-38)} اور درمیانی نقطہ کیا ہے؟}
\usqd{کیمسٹری کے طلبہ کے ایک گروپ نے کیمیائی لیبارٹری میں درج ذیل نمبر حاصل کیے:}{\lr{(29-33)} اور \lr{(44-48)} کے جماعتی تعدد کیا ہے؟}
\usqd{کیمسٹری کے طلبہ کے ایک گروپ نے کیمیائی لیبارٹری میں درج ذیل نمبر حاصل کیے:}{مندرجہ بالا تعددی تقسیم میں جماعتی حدود کی جسامت کیا ہے؟}
\usqd{کیمسٹری کے طلبہ کے ایک گروپ نے کیمیائی لیبارٹری میں درج ذیل نمبر حاصل کیے:}{کس جماعتی گروپ میں طلبہ کی تعداد سب سے کم ہے؟}
\usqd{کیمسٹری کے طلبہ کے ایک گروپ نے کیمیائی لیبارٹری میں درج ذیل نمبر حاصل کیے:}{اس جماعتی حد کی زیریں جماعتی حد کیا ہے جس کا جماعتی تعدد \lr{15} ہے؟}
\usqd{کیمسٹری کے طلبہ کے ایک گروپ نے کیمیائی لیبارٹری میں درج ذیل نمبر حاصل کیے:}{ان طلبہ کی تعداد بتائیں جنھوں نے \lr{24} اور \lr{43} کے درمیان نمبر حاصل کیے۔}
\usq{درج ذیل میں \lr{50} سکوں کو اچھالنے سے کسی تجربے میں ہیڈ آنے کی تعداد دی گئی ہے۔ ان معلومات سے ایک غیر مسلسل تعددی تقسیم کا جدول بنائیں:}
\par\smallskip{\raggedleft\noindent\hspace*{7em}\lr{3 3 4 0 5 4 3 3 1 2 4 5 0 3 2 4 4 0 0 0 5 5}\par\hspace*{7em}\lr{4 3 2 5 3 2 1 3 5 4 3 2 1 3 2 1 3 1 3 1 4 3}\par}
\usq{درج ذیل کا حسابی اوسط معلوم کریں: \lr{4,6,10,12,15,20,25,28,30}}
\usq{درج ذیل کا حسابی اوسط معلوم کریں: \lr{12,18,19,0,-19,-18,-12}}
\usq{درج ذیل کا حسابی اوسط معلوم کریں: \lr{6.5,11,12.3,9,8.1,16,18,20.5,25.}}
\usq{درج ذیل کا حسابی اوسط معلوم کریں: \lr{8,10,12,14,16,20,22.}}
\usq{درج ذیل میں \lr{12} طلبہ کے قد (انچوں میں) ہیں۔ وسطانی قد معلوم کریں: \lr{55,53,54,58,60,61,62,56,57,52,51,63.} مواد کا وسطانیہ معلوم کریں۔}
\usq{درج ذیل میں \lr{10} مزدوروں کی اجرتیں (روپوں میں) ہیں۔ \lr{88,70,72,125,115,95,81,90,95,90} حسابی اوسط معلوم کریں۔}
\usq{درج ذیل میں \lr{10} مزدوروں کی اجرتیں (روپوں میں) ہیں۔ \lr{88,70,72,125,115,95,81,90,95,90} وسطانیہ معلوم کریں۔}
\usq{درج ذیل میں \lr{10} مزدوروں کی اجرتیں (روپوں میں) ہیں۔ \lr{88,70,72,125,115,95,81,90,95,90} عادہ معلوم کریں۔}
\usq{\lr{10} نوجوان لڑکے ایک پیٹرول پمپ سٹیشن پر کام کرتے ہیں۔ ان کی ہفتہ وار اجرتیں درج ذیل ہیں: \lr{4250,4350,4400,4250,4350,4410,4500,4300,4500,4390.} اجرتوں کا حسابی اوسط مختصر کلیے سے، وسطانیہ اور عادہ معلوم کریں۔}
\usq{\lr{45} مدات کا حسابی اوسط \lr{80} ہے۔ ان کا مجموعہ معلوم کریں۔}
\usq{پانچ مدات اس طرح ہیں: \lr{1,4,0,7,9.} ان کا حسابی اوسط، وسطانیہ اور عادہ معلوم کریں۔}
\usq{مواد کے ایک سیٹ میں درج ذیل قیمتیں ہیں: \lr{148,145,160,157,156,160.} ثابت کریں کہ حسابی اوسط \lr{$>$} وسطانیہ \lr{$>$} عادہ}
\usq{درج ذیل معلومات سے حسابی اوسط معلوم کریں: اگر \lr{$D=X-140,\ \sum D=500$} اور \lr{$n=10$} ہو}
\usq{درج ذیل معلومات سے حسابی اوسط معلوم کریں: اگر \lr{$U=\dfrac{x-130}{6},\ \sum U=-150$} اور \lr{$n=15$} ہو}
\usq{درج ذیل معلومات سے حسابی اوسط معلوم کریں: اگر \lr{$D=x-25,\ \sum fD=300$} اور \lr{$\sum f=20$}}
\usq{درج ذیل معلومات سے حسابی اوسط معلوم کریں: اگر \lr{$U=\dfrac{x-120}{5},\ \sum fU=60$} اور \lr{$\sum f=100$}}
\usq{نیچے دیے گئے مواد کا اوزانی وسط معلوم کریں:}
\utabA
\usq{ایک کمپنی اپنے اگلے سال کے مارکیٹنگ بجٹ کی منصوبہ بندی پانچ سالوں میں کر رہی ہے۔ سالانہ بجٹ (ملین میں) \lr{5,7,8,6,7} ہیں۔ اگلے سال کا اوسط بجٹ معلوم کریں۔}
\usq{احمد نے ایک مخصوص امتحان میں درج ذیل نمبر حاصل کیے۔}
\utabB
\usc{اگر مضامین کو بالترتیب \lr{5,4,2,3,2,4} وزن دیا گیا ہو تو ان کا وزانی وسط معلوم کریں۔}
\usq{ایک طالب علم کے پانچ امتحانات میں نمبر \lr{64,75,81,87,90} تھے۔ نمبروں کا حسابی اوسط معلوم کریں۔}
\usq{ایک بلے باز نے درج ذیل دوڑیں بنائیں \lr{40,45,50,52,50,60,56,70.} دیے ہوئے مواد کا حسابی اوسط مختصر کلیہ کی مدد سے معلوم کریں۔}
\usq{کرکٹ کے ایک کھلاڑی نے درج ذیل دوڑیں بنائیں: \lr{8,12,18,13,16,5,50.} دوڑوں کا وسطانیہ معلوم کریں۔}
\usq{\lr{10} طلبہ نے انگریزی میں (\lr{100} میں سے) درج ذیل نمبر حاصل کیے۔ \lr{23,15,35,48,41,5,8,9,11,51.} اس مواد کا وسطانیہ معلوم کریں۔}
\usq{درج ذیل مواد ایک طالب علم کے مختلف مضامین میں نمبروں کی وضاحت کرتا ہے اور ان مضامین کے تفویض کردہ اوزان بھی دیے گئے ہیں:}
\utabC
\usc{اس کا وزانی وسط معلوم کریں۔}
\usq{مندرجہ ذیل کی وضاحت کریں: تعددی تقسیم}
\usq{مندرجہ ذیل کی وضاحت کریں: کالمی نقشہ (غیر مساوی جماعتی حدود کے لیے)}
\usq{مندرجہ ذیل کی وضاحت کریں: حسابی اوسط}
\usq{مندرجہ ذیل کی وضاحت کریں: وسطانیہ}
\usq{45 طلبہ کے بائیولوجی کے ماہانہ ٹیسٹ میں درج ذیل نمبر حاصل کیے:}\ubioB\usc{اوپر دی ہوئی تفصیل سے درج ذیل کے بارے میں بتائیں: پانچویں جماعت کی بالائی حقیقی جماعتی حد}
\usq{45 طلبہ کے بائیولوجی کے ماہانہ ٹیسٹ میں درج ذیل نمبر حاصل کیے:}\ubioA\usc{اوپر دی ہوئی تفصیل سے درج ذیل کے بارے میں بتائیں: تمام جماعتوں کی زیریں حقیقی جماعتی حدود}
\usq{45 طلبہ کے بائیولوجی کے ماہانہ ٹیسٹ میں درج ذیل نمبر حاصل کیے:}\ubioB\usc{اوپر دی ہوئی تفصیل سے درج ذیل کے بارے میں بتائیں: تمام جماعتوں کا درمیانی نقطہ}
\usq{45 طلبہ کے بائیولوجی کے ماہانہ ٹیسٹ میں درج ذیل نمبر حاصل کیے:}\ubioA\usc{اوپر دی ہوئی تفصیل سے درج ذیل کے بارے میں بتائیں: سب سے کم تعداد والی جماعت کی جماعتی حدود}
\usq{سعد کے چھ ٹیسٹوں میں حاصل کردہ نمبر \lr{84,91,72,68,87,78} تھے۔ اس کے نمبروں کا حسابی اوسط معلوم کریں۔}
\usq{عثمان نے آٹھ مرتبہ دو ڈائس پھینکے۔ ہر دفعہ ان کے مجموعہ کو ریکارڈ کر لیا گیا۔ \lr{8,5,6,6,9,4,3,11}۔ ان کے مجموعہ کا وسطانیہ اور عادہ معلوم کریں۔}
\usq{ارشد ایک ڈائس پھینکتا ہے۔ جس کے اطراف میں \lr{L, M, NO, P, U} درج ہوں۔ ڈائس پر صحیح حرف آنے کا احتمال کیا ہے؟}
\usq{شازیہ نے ڈائس کا ایک جوڑا پھینکا۔ احتمال کیا ہو گا جب کہ: نقاط کا مجموعہ کم از کم \lr{4} ہو۔}
\usq{شازیہ نے ڈائس کا ایک جوڑا پھینکا۔ احتمال کیا ہو گا جب کہ: دونوں نقاط کے درمیان کا حاصل ضرب \lr{5} سے \lr{10} کے درمیان ہو۔}
\usq{شازیہ نے ڈائس کا ایک جوڑا پھینکا۔ احتمال کیا ہو گا جب کہ: دونوں نقاط کے درمیان کا فرق \lr{4} کے برابر ہو۔}
\usq{شازیہ نے ڈائس کا ایک جوڑا پھینکا۔ احتمال کیا ہو گا جب کہ: پہلے ڈائس پر کم از کم \lr{5} اور دوسرے ڈائس پر کم از کم \lr{4} ہو۔}
\usq{\lr{“MATHEMATICS”} کے لفظ سے بے ترتیب ایک حرف تہجی منتخب کیا گیا ہے۔ احتمال معلوم کریں جب کہ وہ: حرف علت ہو}
\usq{\lr{“MATHEMATICS”} کے لفظ سے بے ترتیب ایک حرف تہجی منتخب کیا گیا ہے۔ احتمال معلوم کریں جب کہ وہ: حرف صحیح ہو}
\usq{\lr{“MATHEMATICS”} کے لفظ سے بے ترتیب ایک حرف تہجی منتخب کیا گیا ہے۔ احتمال معلوم کریں جب کہ وہ: حرف \lr{E} ہو}
\usq{\lr{“MATHEMATICS”} کے لفظ سے بے ترتیب ایک حرف تہجی منتخب کیا گیا ہے۔ احتمال معلوم کریں جب کہ وہ: حرف \lr{A} ہو}
\usq{\lr{“MATHEMATICS”} کے لفظ سے بے ترتیب ایک حرف تہجی منتخب کیا گیا ہے۔ احتمال معلوم کریں جب کہ وہ: حرف \lr{M} نہ ہو}
\usq{\lr{“MATHEMATICS”} کے لفظ سے بے ترتیب ایک حرف تہجی منتخب کیا گیا ہے۔ احتمال معلوم کریں جب کہ وہ: حرف \lr{T} نہ ہو}
\usq{احتمال کیا ہے؟}
\usq{احتمال کا فارمولا لکھیں۔}
\usq{تجربہ کی تعریف کریں۔}
\usq{نتائج کی تعریف کریں۔}
\usq{موافق نتیجہ کی تعریف کریں۔}
\usq{سیمپل اسپیس کی تعریف کریں۔}
\usq{واقعہ کی تعریف کریں۔}
\usq{سیمپل پوائنٹ کی تعریف کریں۔}
\usq{یقینی واقعہ کیا ہے؟}
\usq{ناممکن واقعہ کیا ہے؟}
\usq{ممکنہ واقعہ کیا ہے؟}
\usq{غیر ممکنہ واقعہ کیا ہے؟}
\usq{مساوی ممکنہ واقعات کیا ہوتے ہیں؟}
\usq{عبدالرحیم ایک ڈائس پھینکتا ہے۔ \lr{3} پر تقسیم ہونے والے عدد کا احتمال کیا ہوتا ہے؟}
\usq{اگر ذیشان \lr{2} ڈائس پھینکتا ہے تو احتمال معلوم کریں کہ دونوں ڈائس پر جفت اعداد ہوں۔}
\usq{اگر ذیشان \lr{2} ڈائس پھینکتا ہے تو احتمال معلوم کریں کہ دونوں ڈائس پر \lr{3} کے اضعاف ہوں}
\usq{اگر ذیشان \lr{2} ڈائس پھینکتا ہے تو احتمال معلوم کریں کہ پہلے ڈائس پر جفت عدد اور دوسرے ڈائس پر ہندسہ \lr{“3”} ہو}
\usq{اگر ذیشان \lr{2} ڈائس پھینکتا ہے تو احتمال معلوم کریں کہ پہلے ڈائس پر ہندسہ کم از کم \lr{3} ہو اور دوسرے ڈائس پر ہندسہ \lr{4} ہو۔}
\usq{زبیر ایک ڈائس پھینکتا ہے۔ ہندسہ \lr{6} حاصل نا ہونے کا احتمال کیا ہو گا؟}
\usq{اگر دو ڈائس پھینکے جائیں تو احتمال کیا ہو گا کہ: دونوں ڈائس کا مجموعہ \lr{8} نا ہو}
\usq{اگر دو ڈائس پھینکے جائیں تو احتمال کیا ہو گا کہ: دو مرتبہ \lr{6} نا آئے}
\usq{تاش کے \lr{52} پتوں (کارڈز) میں سے ایک پتہ بے ترتیب منتخب کیا گیا۔ ایک پتہ منتخب کرنے کا احتمال کیا ہے اگر: پان کا پتہ ہو}
\usq{تاش کے \lr{52} پتوں (کارڈز) میں سے ایک پتہ بے ترتیب منتخب کیا گیا۔ ایک پتہ منتخب کرنے کا احتمال کیا ہے اگر: نا حکم نا ہی پان کا پتہ ہو}
\usq{دو ڈائس کو \lr{500} دفعہ پھینکا گیا، اگر مجموعہ \lr{8} سے زیادہ ہو تو متوقع تعدد کیا ہو گی؟}
\usq{اگر ایک شخص تین سکوں کو ایک دفعہ اچھالے اور کم از کم \lr{2} ہیڈ حاصل کرے تو اسے \lr{120} روپے انعام ملنے کی کیا توقع ہو گی؟}
\usq{نسبتی تعدد کیا ہے؟}
\usq{نسبتی تعدد کیسے معلوم کیا جاتا ہے؟}
\usq{متوقع تعدد کی تعریف کریں۔}
\usq{ایک ڈائس \lr{300} مرتبہ اچھالا گیا۔ ہندسہ \lr{1} یا \lr{6} حاصل ہونے کی اوسط تعداد معلوم کریں۔}
\usq{اگر خراب بولٹ کا احتمال \lr{0.3} ہو تو کل \lr{800} بولٹ میں سے غیر ناقص بولٹ کی تعداد معلوم کریں۔}

\end{document}
